\documentclass[11pt]{article}
\usepackage[margin=1in]{geometry}
\usepackage{graphicx}
\usepackage{amsmath,amssymb}
\usepackage{booktabs}
\usepackage{caption}
\usepackage{subcaption}
\usepackage{float}
\usepackage{array}
\usepackage{siunitx}
\usepackage[hidelinks]{hyperref}
\graphicspath{{./}{figs/}{figures/}{pub_figs/}{openfoam/doc_figs/}{openfoam/chamber42_3d/calibration/}}
\captionsetup{font=small,labelfont=bf}
\setlength{\parskip}{4pt}

\title{\textbf{A Subcooled Pool-Boiling Chamber with a Submerged Tube-Bundle Condenser:\\
A Validated Reduced-Order Model and a Three-Dimensional VOF Phase-Change Framework}}
\author{[Author list]\\
\small Thermal Analysis Laboratory, Department of Mechanical Engineering\\
\small Rochester Institute of Technology, Rochester, NY, USA}
\date{\today}

\begin{document}
\maketitle

\begin{abstract}
\noindent
This document combines two complementary models of a closed, subcooled pool-boiling
chamber in which a heated chip boils a working fluid and the vapor is rejected to coolant
through a partially submerged tube-bundle condenser, a configuration developed for
high-heat-flux data-center cooling. \textbf{Part I} presents a validated reduced-order
predictive thermal model: a mechanistically grounded resistance network with a small set of
physically interpretable constants that predicts chip surface temperature versus heat flux
across two evaporator surfaces (plain copper and open microchannel), two condenser bundles
(33-tube and 42-tube), and four coolant inlet temperatures, reproducing 261 experimental
points with a root-mean-square error of \SI{4.3}{K}. \textbf{Part II} presents a
three-dimensional volume-of-fluid phase-change simulation framework for the same apparatus,
built on the true CAD geometry: a verified body-fitted mesh, the governing equations and
numerics, the condenser coolant-side flow routing, and a calibration-to-validation workflow
in which the measured boiling curves and thermal resistances of Part~I are the validation
targets. Together they span the modeling spectrum from a fast, validated predictor to a
high-fidelity framework, under a consistent standard that distinguishes validated results,
in-sample fits, and pending CFD output.
\end{abstract}

\tableofcontents
\clearpage

\part{Validated Reduced-Order Predictive Model}

\section{Introduction}
Two-phase immersion cooling of high-power electronics relies on nucleate boiling at a
heated surface and condensation of the generated vapor at a heat exchanger. In a closed
chamber the two processes are coupled: the chamber pressure, hence the saturation
temperature $T_{sat}$ at which the fluid boils and condenses, settles so that the
condenser rejects exactly the heat the chip dissipates. Predicting chip temperature for
an untested geometry therefore requires not only a boiling correlation but also a
prediction of the operating point itself. The objective is a genuinely predictive model
that forecasts chip surface temperature for new condenser and evaporator geometries,
coolant inlet temperatures, and chip surfaces, rather than curve-fitting a single data
set. The model is a reduced-order ``gray box'': mechanistic correlations carry the
physics, and a minimal number of physically interpretable calibrated constants absorb
what cannot be resolved from geometry alone.

\section{Literature review}

\paragraph{Thermal management context.}
The rising power density of microprocessors and the growth of data-center workloads have
made conventional air cooling insufficient, motivating direct liquid and two-phase
schemes \cite{mudawar2001,azarifar2024}. Two-phase immersion cooling, in which a chip is
immersed in a dielectric or water bath that boils on contact and condenses elsewhere in
the enclosure, offers very high heat transfer coefficients and near-isothermal operation
\cite{birbarah2020,wu2025}. The present chamber is a closed realization of this concept
with a submerged condenser, so its performance depends jointly on the boiling surface and
the condenser.

\paragraph{The present apparatus.}
The chamber modeled here was introduced by Shukla and Kandlikar \cite{shukla2024}, who combined
subcooled pool boiling with a partially submerged tube-bundle condenser in a compact form factor
that fits a 1U/2U server rack. Partly submerging the condenser augments subcooled boiling and adds
direct-contact condensation, and high-speed imaging confirmed submerged condensation with small,
discrete vapor bubbles. The reported thermal test vehicle dissipated about $400~\mathrm{W}$ at a
surface temperature below $80~^\circ$C on Novec-7000, and over $1~\mathrm{kW}$ from a
$34.5\times32~\mathrm{mm}$ plain copper chip on water, establishing the operating envelope this
model is built to predict.

\paragraph{Pool boiling and the boiling curve.}
The characteristic boiling curve, with its single-phase, nucleate, transition and film
regimes, was established by Nukiyama \cite{nukiyama1934} and treated comprehensively in
the monographs of Carey \cite{carey2008} and Incropera et al.\ \cite{incropera2007}. The
most widely used nucleate-boiling correlation is that of Rohsenow \cite{rohsenow1952}, in
which a surface--fluid coefficient $C_{sf}$ encapsulates the combined effect of surface
finish, wettability and fluid. Mikic and Rohsenow \cite{mikic1969} refined the treatment
of surface characteristics, and Pioro \cite{pioro1999} compiled experimental $C_{sf}$
values, emphasizing that the coefficient is a property of the specific surface--fluid pair
and must be calibrated. This motivates fitting $C_{sf}$ once per chip surface and treating
it as a transferable constant.

\paragraph{Critical heat flux.}
The peak nucleate flux, or critical heat flux (CHF), was first described by hydrodynamic
instability theory \cite{zuber1959,kutateladze1948}, later extended to finite and oriented
surfaces by Lienhard and Dhir \cite{lienhard1973}. Kandlikar \cite{kandlikar2001}
introduced a force-balance model incorporating the dynamic receding contact angle and
surface orientation; Mudawar and co-workers developed complementary near-saturated CHF
models \cite{mudawar1997}. Because no burnout data beyond the measured CHF points were
available, CHF is anchored to measurement rather than modeled mechanistically.

\paragraph{Enhanced surfaces and open microchannels.}
Surface structuring is the principal route to higher CHF and heat transfer coefficient.
Cooke and Kandlikar \cite{cooke2011,cooke2012} showed that open microchannels milled into
copper dissipate fluxes on the order of $240~\mathrm{W/cm^2}$ at low wall superheat, and
identified channel rewetting and separate liquid--vapor pathways, rather than mere area
increase, as the dominant mechanism. Subsequent work combined microchannels with sintered
or electrodeposited porous coatings \cite{jaikumar2015,patil2014} and demonstrated
bubble-induced convective liquid feeding \cite{jaikumar2016}. The chip modeled here uses
the Cooke--Kandlikar geometry, represented by a fixed geometric area augmentation with a
reduced $C_{sf}$.

\paragraph{Dielectric working fluids.}
Immersion cooling favors dielectric coolants, fluorocarbons and hydrofluoroethers such as
Novec/HFE-7000, for their electrical inertness and low, near-ambient saturation temperature
(about $34~^\circ$C for HFE-7000 at atmospheric pressure), at the cost of a much lower critical heat
flux than water. El-Genk and Pourghasemi \cite{elgenk2021a,elgenk2021b} measured saturated and
subcooled boiling of HFE-7000 on copper of varying roughness and orientation, correlating both the
heat transfer coefficient and CHF against surface roughness, subcooling, and inclination; they
report CHF on plain copper near $20$--$30~\mathrm{W/cm^2}$, rising $50$--$70\%$ at $30~\mathrm{K}$
subcooling. These dependences, roughness, subcooling, and contact angle, are the levers the
present CHF closure must capture when the model is extended from water to HFE-7000.

\paragraph{Condensation.}
The condenser side rests on the classical film-condensation analysis of Nusselt
\cite{nusselt1916}, extended to dropwise condensation \cite{rose2002} and to tube bundles
in the standard references \cite{collier1994}. A central finding here is that, for this
geometry, condensation is not the limiting resistance.

\paragraph{Internal convection.}
The coolant-side resistance is governed by internal forced convection. For the short,
small-bore tubes used here the flow is laminar and thermally developing, a regime
described by Hausen's entrance-region correlation \cite{hausen1943} and tabulated by Shah
and London \cite{shah1978}; the turbulent limit is covered by Gnielinski
\cite{gnielinski1976}, and external natural convection by Churchill and Chu
\cite{churchill1975}. Recognizing the developing-laminar character of the coolant flow is
what corrects an earlier over-prediction of the operating point.

\paragraph{System-level studies and reduced-order models.}
At the system level, immersion cooling is pursued in both single-phase \cite{shrigondekar2023} and
two-phase \cite{wu2025} forms, and recent multi-chip and loop studies emphasize that overall
performance is set as much by the condenser and the coolant loop as by the boiling surface. Detailed
CFD of such chambers is costly, which has motivated semi-empirical and lumped models for design-space
exploration; the gray-box developed here follows that line, pairing mechanistic correlations with a
small set of calibrated, physically interpretable constants rather than a full numerical solution.

\paragraph{Closed-chamber coupling and the gap addressed here.}
Closed two-phase enclosures and thermosyphons couple evaporation and condensation through
the vapor space \cite{faghri1995}, and non-condensable gases and headspace pressure shift
the operating point. Yet most boiling and condensation studies treat the two processes in
isolation, and few predict the coupled operating point of a closed chamber from geometry
alone. The present model addresses that gap with a calibrated, mechanistically grounded
operating-point sub-model.

\section{Methods}

\subsection{Chip and condenser geometry}
Both chips share an $11.04~\mathrm{cm^2}$ ($34.5\times32$~mm) copper footprint and are
mounted at the base of the chamber. The plain chip is polished to $R_a\approx2.5~\mu$m.
The enhanced chip carries open microchannels of the Cooke--Kandlikar type
\cite{cooke2012}: rectangular channels $381~\mu$m wide and $400~\mu$m deep separated by
$250~\mu$m fins, increasing the wetted area and providing rewetting pathways. Three
submerged copper condensers are considered (Table~\ref{tab:geom}). The 33-tube design uses
relatively large $4.76$~mm tubes with external (open) manifolds; the 42-tube design uses
denser, smaller $3.175$~mm tubes with in-chamber manifolds, placing a lower, more
submerged bank in the pool; the 66-tube design exists in CAD but has no experimental data.
The coolant is water, supplied at $\sim\!78~\mathrm{g/s}$ (33-tube) and
$\sim\!36~\mathrm{g/s}$ (42-tube).

\begin{table}[H]
\centering
\caption{Condenser and chip geometry.}
\label{tab:geom}
\begin{tabular}{lccccc}
\toprule
Condenser & Tubes $N$ & OD (mm) & ID (mm) & $L$ (mm) & Ext.\ area (cm$^2$)\\
\midrule
33-tube (open manifolds)       & 33 & 4.76  & 3.14 & 95.0 & 469 \\
42-tube (in-chamber manifolds) & 42 & 3.175 & 1.39 & 62.6 & 262 \\
66-tube (no data yet)          & 66 & 2.00  & 1.80 & 95.0 & 394 \\
\midrule
\multicolumn{6}{l}{\textit{Chip surfaces}}\\
Plain        & \multicolumn{5}{l}{polished copper, $R_a\approx2.5~\mu$m}\\
Microchannel & \multicolumn{5}{l}{open channels $381~\mu$m wide $\times\,400~\mu$m deep, $250~\mu$m fins}\\
\bottomrule
\end{tabular}
\end{table}

The two condenser bundles are shown in cross-section in Figure~\ref{fig:xsec}, reconstructed
from the CAD. The 33-tube bundle is a single bank of 33 relatively large straight tubes in five
staggered rows; the 42-tube bundle is denser, with 42 smaller-diameter straight tubes in four
staggered rows fed by in-chamber manifolds. All tube axes are parallel; the tubes are straight,
not U-shaped.

\begin{figure}[H]
\centering
\includegraphics[width=0.74\textwidth]{mc_xsec.png}
\caption{Condenser cross-sections drawn to a common scale from the CAD geometry, looking along
the tube axis (coolant into the page). The 33-tube bundle uses larger straight tubes at wider
pitch; the 42-tube bundle uses more, smaller straight tubes at tighter pitch.}
\label{fig:xsec}
\end{figure}

\subsection{Modeling domain}
The model resolves the chip-to-coolant path as the reduced thermal network of
Figure~\ref{fig:domain}. Three temperatures are tracked: the chip surface $T_{surf}$, the
bulk (pool) liquid $T_{bulk}$, and the coolant inlet $T_{in}$. Heat $Q$ flows from the
chip into the pool across a boiling/chip resistance and from the pool to the coolant across
a condenser resistance, the two acting in series. The pool is treated as a single
well-mixed reservoir at $T_{bulk}=T_{sat}-\Delta T_{sub}$; the vapor space and chamber
pressure are represented implicitly through the saturation temperature; and the condenser
is reduced to a coolant-side-limited heat exchanger characterized by the operating-point
map. The chip surface is a lumped boiling node rather than a spatially resolved field. This
reduction is justified a posteriori by the resistance breakdown, which shows the path is
dominated by two clearly separated resistances.

\begin{figure}[H]
\centering
\includegraphics[width=\textwidth]{mc_domain.png}
\caption{Physical domain (left) and the reduced thermal network (right). The chip boils the
subcooled pool; vapor condenses on the submerged bundle and rejects heat to the internal
water coolant. The chip-to-coolant path reduces to two series resistances about the bulk
liquid temperature.}
\label{fig:domain}
\end{figure}

\subsection{Governing equations}
The chip dissipates heat by simultaneous single-phase convection into the subcooled pool
and nucleate boiling. The nucleate component follows the Rohsenow correlation
\cite{rohsenow1952},
\begin{equation}
q''_{nb}(T_s) = \mu_l h_{fg}\sqrt{\frac{g(\rho_l-\rho_v)}{\sigma}}
\left[\frac{c_{p,l}\,(T_s-T_{sat})}{C_{sf}\,h_{fg}\,Pr_l^{\,s}}\right]^{3},
\qquad T_s>T_{sat},
\label{eq:rohsenow}
\end{equation}
with $s=1$ for water. The single-phase component is a horizontal-surface natural-convection
coefficient scaled by a per-chamber factor $F_{nc}$,
\begin{equation}
q''_{nc}(T_s) = F_{nc}\, h_{nc}(T_s,T_{bulk})\,(T_s-T_{bulk}).
\label{eq:nc}
\end{equation}
The mechanisms combine in quadrature, and the microchannel surface uses the same
expressions multiplied by a geometric area-augmentation ratio $A_r$:
\begin{equation}
q'' = A_r\sqrt{q''_{nc}(T_s)^2 + q''_{nb}(T_s)^2}.
\label{eq:blend}
\end{equation}
The microchannel critical heat flux exceeds the plain value by a near-constant,
geometry-independent increment,
\begin{equation}
q''_{CHF,\,micro} = q''_{CHF,\,plain} + \Delta q''_{CHF},
\qquad \Delta q''_{CHF}\approx 60~\mathrm{W/cm^2}.
\label{eq:chf}
\end{equation}
The chamber operating point closes the system. A resistance breakdown
(Section~\ref{sec:results}) shows the condenser is coolant-side limited, with internal flow
developing-laminar ($Re\!\sim\!1100$) and film condensation nearly ideal and almost
negligible; the coolant-side conductance improves as the system warms, so $T_{sat}$
responds to coolant inlet with a gain below unity. The developing-laminar Nusselt number
\cite{hausen1943},
\begin{equation}
Nu = 3.66 + \frac{0.0668\,Gz}{1+0.04\,Gz^{2/3}},\qquad Gz=\frac{D_i}{L}\,Re\,Pr,
\label{eq:hausen}
\end{equation}
reproduces the magnitude and temperature dependence of the dominant resistance and motivates
a calibrated linear operating-point map per chamber,
\begin{align}
T_{sat} &= a\,T_{in} + b\,Q + c, \label{eq:tsat}\\
\Delta T_{sub} &= \max\!\big(d\,T_{in} + e\,Q + f,\;0\big),\label{eq:sub}\\
Q &= q''\,A_{\mathrm{foot}}, \qquad A_{\mathrm{foot}}=11.04~\mathrm{cm^2}.\nonumber
\end{align}
The chip-to-coolant path is split at the bulk liquid into two series resistances,
\begin{equation}
R_{chip}=\frac{T_s-T_{bulk}}{Q},\qquad R_{cond}=\frac{T_{bulk}-T_{in}}{Q},\qquad
R_{tot}=R_{chip}+R_{cond}=\frac{T_s-T_{in}}{Q}.
\label{eq:resistance}
\end{equation}

\subsection{Solver setup}
The model is implemented in Python (NumPy/SciPy). Thermophysical properties are evaluated
with CoolProp for water (\texttt{Water}). HFE-7000 is not available in the CoolProp build used
here, so its properties are taken from the literature \cite{elgenk2021a} and held fixed
($h_{fg}=142~\mathrm{kJ/kg}$, $\sigma=12.4~\mathrm{mN/m}$, $T_{sat}\approx34~^\circ$C at $1$~atm).
Properties are evaluated at the saturation or film temperature and cached. The fully predictive chain proceeds as follows: from the
inputs $(q'',T_{in},\text{geometry})$ the operating-point map, Eqs.~\eqref{eq:tsat}--\eqref{eq:sub},
returns $T_{sat}$ and $\Delta T_{sub}$, hence $T_{bulk}$; the chip relation,
Eq.~\eqref{eq:blend}, is then solved for the surface temperature $T_s$ by Brent's method on
the bracket $[\,T_{bulk},\,T_{sat}+260\,\mathrm{K}\,]$, widened to $+1000$~K (water) or $+200$~K
(HFE-7000) if the root is not initially bracketed, to a tolerance of $10^{-6}~\mathrm{K}$. No measured operating condition enters the prediction. Boiling curves
and parity plots are generated by sweeping $q''$ from $2~\mathrm{W/cm^2}$ to the
configuration CHF. The model contains no spatial PDE; it is a reduced-order network of
algebraic closures, so the only numerical parameters are the root-finder tolerance, the
flux-sweep resolution, and the discretization of the mechanistic coolant-side estimate.

\subsection{Discretization sensitivity}
Two numerical resolutions were verified (Figure~\ref{fig:mesh}). First, the mechanistic
coolant-side estimate that motivates the resistance coefficient $b$ integrates the local
developing-laminar Nusselt number along each tube by a one-dimensional finite-volume
marching scheme: the tube is divided into $N$ axial segments, the coolant temperature is
advanced segment by segment under an isothermal condensing wall, and the length-averaged
internal Nusselt number is recovered from the log-mean temperature difference. The
length-averaged $\overline{Nu}$ converges monotonically, changing by less than $0.4\%$
between $N=128$ and $N=512$ for both chambers. Second, the predicted boiling curves and
parity data are sampled on a flux sweep; the curve is grid-independent to better than
$0.05~\mathrm{K}$ beyond $80$ points. Production runs use $256$ axial segments and $90$
flux points, and the nonlinear surface-temperature solve converges to below
$10^{-4}~\mathrm{K}$, so all reported results are numerically converged.

\begin{figure}[H]
\centering
\includegraphics[width=\textwidth]{mc_mesh.png}
\caption{Discretization sensitivity. Left: length-averaged coolant-side Nusselt number versus
axial segments per tube from the 1D marching scheme; both chambers plateau by a few hundred
segments. Right: maximum deviation of the predicted boiling curve from a fine reference versus
flux-sweep resolution.}
\label{fig:mesh}
\end{figure}

\subsection{Coefficient determination and justification}
\label{sec:coeff}
The calibrated constants are listed in Table~\ref{tab:params}; how each was obtained and why
it is defensible follows.

\textit{Surface--fluid coefficients $C_{sf}$.} Fit by least squares of
Eq.~\eqref{eq:rohsenow} to the measured wall-superheat--flux data on each chip, all coolant
inlets pooled. Both calibrated values land on tabulated Rohsenow coefficients for the matching
copper finish (Table~\ref{tab:csf}). The plain value $0.0131$ agrees with water on polished
copper, $0.0128$--$0.0130$ \cite{rohsenow1952,vachon1968}, to within $2\%$. The microchannel
value $0.0067$, roughly half, coincides with the classic value for water on \emph{scored} copper,
$0.0068$ \cite{vachon1968}, which is physically apt: the open microchannels are an aggressive
grooving of the surface that raises the active nucleation-site density \cite{cooke2012}. That two
coefficients fit only to this chamber's data reproduce independent handbook values for the
corresponding surface finish is direct evidence they are genuine surface--fluid properties rather
than tuning constants. Section~\ref{sec:crossval} adds a cross-chamber transfer test, and the
identical plain boiling curve on both chambers confirms chamber independence.

\begin{table}[H]
\centering
\caption{Calibrated $C_{sf}$ against tabulated Rohsenow values for water on copper, ordered by
value to pair each chip with its matching finish. Handbook values from
\cite{rohsenow1952,vachon1968}.}
\label{tab:csf}
\begin{tabular}{lc}
\toprule
Surface (water on copper) & $C_{sf}$\\
\midrule
Scored copper (handbook) & 0.0068\\
\textbf{Microchannel chip (this work)} & \textbf{0.0067}\\
Polished copper (handbook) & 0.0128--0.0130\\
\textbf{Plain chip (this work)} & \textbf{0.0131}\\
Lapped copper (handbook) & 0.0147\\
\bottomrule
\end{tabular}
\end{table}

\textit{Area augmentation $A_r=2.17$.} Not fitted. It is the geometric wetted-to-footprint
area ratio of the Cooke--Kandlikar microchannel, computed from the channel and fin dimensions.
Using the geometric value rather than a free multiplier preserves the predictive intent.

\textit{CHF increment $\Delta q''_{CHF}=+60~\mathrm{W/cm^2}$.} Obtained from the measured
plain-to-microchannel CHF differences, which are $+64$ and $+59~\mathrm{W/cm^2}$ on the two
chambers. The additive constant fits both far better than multiplying by $A_r$, which would
over-predict the microchannel CHF to $\sim\!247$ and $\sim\!141~\mathrm{W/cm^2}$. A roughly
fixed flux increment is consistent with a rewetting-limited enhancement that is largely
independent of bundle geometry.

\textit{Single-phase factor $F_{nc}$.} Fit per chamber on the low-flux, pre-boiling data to
match the single-phase slope. It lumps subcooled-boiling onset and bulk circulation and is
larger for the open 33-tube ($9.0$) than the confined 42-tube ($3.5$), consistent with stronger
circulation in the open layout. It is chamber-specific, not transferable.

\textit{Operating-point coefficients $(a,b,c)$.} Per-chamber multilinear regression of measured
$T_{sat}$ on coolant inlet and chamber heat load, pooling plain and microchannel points across
all coolant inlets. The coolant gain $a$ is below unity ($0.26$ open, $0.71$ confined),
expressing the coolant-side limitation and the conductance improvement as the chamber warms; the
slope $b$ is an effective resistance whose magnitude is consistent with the 1D developing-laminar
estimate of Figure~\ref{fig:mesh}; $c$ is an offset. The regressions give a saturation-temperature
RMSE of about $4~\mathrm{K}$ (33-tube) and $9~\mathrm{K}$ (42-tube).

\textit{Subcooling coefficients $(d,e,f)$.} Per-chamber regression of measured subcooling on
coolant inlet and heat load. The 33-tube fit is well determined ($R^2\approx0.81$); the 42-tube
fit is essentially scatter ($R^2\approx0.06$) and is flagged as a limitation.

\begin{table}[H]
\centering
\caption{Calibrated model constants. $C_{sf}$ and $A_r$ are transferable; $F_{nc}$ and the
operating-point coefficients are per chamber. $b$ and $e$ are in K/W.}
\label{tab:params}
\begin{tabular}{llcc}
\toprule
Group & Constant & 33-tube & 42-tube\\
\midrule
Evaporator & $C_{sf}$ (plain) & \multicolumn{2}{c}{$0.0131$}\\
           & $C_{sf}$ (microchannel) & \multicolumn{2}{c}{$0.0067$}\\
           & Area augmentation $A_r$ (microchannel) & \multicolumn{2}{c}{$2.17$}\\
           & Single-phase factor $F_{nc}$ & $9.0$ & $3.5$\\
\midrule
CHF        & $q''_{CHF}$ plain (W/cm$^2$) & $114$ & $65$\\
           & Microchannel increment $\Delta q''_{CHF}$ (W/cm$^2$) & \multicolumn{2}{c}{$+60$}\\
\midrule
Operating  & $a$ (coolant gain) & $0.262$ & $0.710$\\
$T_{sat}$  & $b$ (K/W) & $5.99\times10^{-3}$ & $2.41\times10^{-2}$\\
           & $c$ (K) & $44.8$ & $33.1$\\
\midrule
Subcooling & $d$ & $-0.639$ & $-0.217$\\
           & $e$ (K/W) & $-6.80\times10^{-3}$ & $-7.01\times10^{-4}$\\
           & $f$ (K) & $40.8$ & $19.6$\\
\bottomrule
\end{tabular}
\end{table}

\subsection{Microchannel coefficient identifiability}
\label{sec:microfit}
Because the open-microchannel chip is the same physical surface in both chambers, its surface
coefficient should be a single transferable value, and refitting it tests that. Fit chip-side
against the measured operating point, the microchannel $C_{sf}$ minimizes at 0.0070 on the
33-tube and 0.0064 on the 42-tube (Fig.~\ref{fig:microfit}a), bracketing the adopted pooled
value 0.0067; the two chambers agree to within the width of a shallow minimum, and refitting
$C_{sf}$ alone changes the chip-side RMSE by less than 0.01~K. The coefficient is therefore well
determined and chamber-independent, consistent with the plain-chip transfer test of
Section~\ref{sec:crossval}.

Freeing the area augmentation $A_r$ jointly with $C_{sf}$ does not help. The two parameters both
multiply the boiling heat transfer, so they trade off along a degenerate valley
(Fig.~\ref{fig:microfit}b): pooled, the minimum drifts to $(0.0084, 3.03)$ for a chip-side gain
of only 0.4~K, while the per-chamber joint fits split to $(0.0064, 1.93)$ and $(0.0091, 3.50)$
and, applied end to end, trade one chamber for the other, improving the 42-tube microchannel RMSE
from 5.2 to 4.5~K but worsening the 33-tube from 4.8 to 5.2~K. Holding $A_r$ at its geometric
value 2.17 and $C_{sf}$ at the pooled 0.0067 is thus the parsimonious choice: it keeps the area a
measured quantity rather than a fitted one and avoids overfitting a two-parameter enhancement to
limited data. The residual microchannel error is not a deficiency of the surface coefficient; with
the operating point taken from the data, the chip-side RMSE (5.6 and 7.1~K) exceeds the end-to-end
RMSE (4.8 and 5.2~K), which means the smooth operating-point map averages out scatter that the raw
measured saturation and subcooling carry, and that scatter, not the chip coefficient, sets the
floor.

\begin{figure}[H]
\centering
\includegraphics[width=\textwidth]{mc_microfit.png}
\caption{Microchannel coefficient identifiability. (a) Chip-side RMSE versus microchannel
$C_{sf}$ with the area augmentation fixed: both chambers minimize near the adopted 0.0067.
(b) Pooled chip-side RMSE over $(C_{sf}, A_r)$: the two parameters trade off along a degenerate
valley, and the per-chamber joint fits (markers) split apart, so the single geometric area plus
the pooled coefficient is the robust choice.}
\label{fig:microfit}
\end{figure}

\subsection{Parameter sensitivity}
\label{sec:sens}
Figure~\ref{fig:sens} shows how far the predicted chip temperature moves when each calibrated
constant is perturbed by $\pm10\%$, averaged over all points. The ranking is informative. The
operating-point offset $c$ dominates ($3.2~\mathrm{K}$) and the coolant gain $a$ is second
($1.5~\mathrm{K}$), because both shift the predicted saturation temperature, which propagates
almost directly into chip temperature. The boiling and single-phase constants are secondary: a
$\pm10\%$ change in $C_{sf}$ moves the chip temperature under $1~\mathrm{K}$ and $A_r$ under
$0.3~\mathrm{K}$. Two conclusions follow. The model is not hypersensitive to any single chip
constant, so the few-percent uncertainty in $C_{sf}$ or $A_r$ established above is not limiting.
And the accuracy is governed mainly by the empirical operating-point map, the same component the
out-of-sample test (Section~\ref{sec:crossval}) identifies as the weakest link and that a direct
condenser closure would most improve.

\begin{figure}[H]
\centering
\includegraphics[width=0.82\textwidth]{mc_sensitivity.png}
\caption{Sensitivity of predicted chip temperature to a $\pm10\%$ change in each calibrated
constant, averaged over all 261 points. Operating-point constants (orange) dominate; chip
constants (blue) are secondary.}
\label{fig:sens}
\end{figure}

\subsection{Measurement uncertainty}
\label{sec:uncertainty}
The error metrics are referenced to the measured chip temperature, which carries its own
experimental uncertainty. Following the laboratory's established procedure
\cite{mody2022,shukla2024}, uncertainties are propagated by the root-sum-square
(Kline--McClintock) method, with contributions from thermocouple calibration, the copper
thermal conductivity, the thermocouple spacing, and the instrument least count. The chip
surface temperature is obtained by extrapolating the embedded thermocouple readings to the
boiling surface through the copper conduction length:
\begin{equation}
T_{surf} = T_{tc} - q''\,\frac{\Delta x}{k_{Cu}}, \qquad
U_{T_{surf}}^2 = U_{T_{tc}}^2 + \Big(\tfrac{\Delta x}{k_{Cu}}\Big)^2 U_{q''}^2
 + \Big(\tfrac{q''}{k_{Cu}}\Big)^2 U_{\Delta x}^2 .
\end{equation}
The calibrated thermocouple uncertainty is $U_{T_{tc}}=\pm0.05~\mathrm{K}$, the
sensor-to-surface spacing $\Delta x = 5~\mathrm{mm}$ at $U_{\Delta x}=\pm0.01~\mathrm{mm}$, and
the copper conductivity $k_{Cu}=401\pm9~\mathrm{W\,m^{-1}K^{-1}}$ ($2.2\%$). The heat flux
follows from the one-dimensional gradient $q''=-k_{Cu}\,(dT/dx)$ fit to the thermocouple column;
its relative uncertainty is dominated by the random scatter in that gradient at low flux and by
the systematic ($k_{Cu}$, spacing) floor at high flux. For the present data it falls from
$10.9\%$ at the lowest flux ($\sim9~\mathrm{W\,cm^{-2}}$) to $4.2\%$ at the peak logged flux
($37~\mathrm{W\,cm^{-2}}$), the high-flux value being bias-limited at about $2.5\%$; the
conduction correction then governs $U_{T_{surf}}$, which stays within $\pm0.2~\mathrm{K}$ across
the range. The saturation temperature is set by the chamber pressure through the vapor-pressure
slope, $U_{T_{sat}}=(dT_{sat}/dP)\,U_P$ with $dT_{sat}/dP\approx0.28~\mathrm{K/kPa}$ for water
and HFE-7000 near their saturation, so a transducer uncertainty of order $0.5~\mathrm{kPa}$ gives
$U_{T_{sat}}\approx\pm0.15~\mathrm{K}$ and a wall-superheat uncertainty within
$\pm0.25~\mathrm{K}$. The coolant-side balance $Q=\dot m\,c_p\,(T_{out}-T_{in})$ provides an
independent check on the chip flux, and the reported critical heat fluxes carry about
$\pm4\%$. These bands set the noise floor against which the model error should be read: at
$4.3~\mathrm{K}$ the chip-temperature RMSE exceeds the $\sim\pm0.2~\mathrm{K}$ measurement
uncertainty by more than an order of magnitude, so the residual reflects model structure,
chiefly the operating-point regression, rather than measurement noise.

\section{Results and validation}
\label{sec:results}

\subsection{Model versus experiment for each chip and condenser}
Figure~\ref{fig:mve} compares the fully predictive model against experiment for each chip and
condenser separately, plotting heat flux against chip surface temperature with the model swept to
CHF. Per-configuration errors are in Table~\ref{tab:rmse}; the aggregate parity is
Figure~\ref{fig:parity}.

\begin{figure}[H]
\centering
\includegraphics[width=0.49\textwidth]{mve_33_plain.png}\hfill
\includegraphics[width=0.49\textwidth]{mve_33_micro.png}\\[3pt]
\includegraphics[width=0.49\textwidth]{mve_42_plain.png}\hfill
\includegraphics[width=0.49\textwidth]{mve_42_micro.png}
\caption{Heat flux versus chip surface temperature for each chip and condenser, model swept to
CHF. Markers are measured (distinct symbol per coolant inlet); lines are the fully predictive
model.}
\label{fig:mve}
\end{figure}

\begin{table}[H]
\centering
\caption{Chip-temperature prediction error (fully predictive).}
\label{tab:rmse}
\begin{tabular}{lcc}
\toprule
Configuration & RMSE (K) & CHF (W/cm$^2$)\\
\midrule
33-tube, plain        & 3.2 & 114\\
42-tube, plain        & 4.5 & 65\\
33-tube, microchannel & 4.9 & 174\\
42-tube, microchannel & 5.2 & 125\\
\midrule
All points            & 4.4 & $-$\\
\bottomrule
\end{tabular}
\end{table}

\begin{figure}[H]
\centering
\includegraphics[width=0.62\textwidth]{mc_endtoend.png}
\caption{End-to-end parity of predicted versus measured chip surface temperature. The shaded band
is $\pm15\%$.}
\label{fig:parity}
\end{figure}

\subsection{Out-of-sample validation}
\label{sec:crossval}
The errors above are computed on the same points used for calibration. To test the central claim
that the model predicts rather than fits, two out-of-sample exercises were run
(Figure~\ref{fig:crossval}, Table~\ref{tab:crossval}). First, a leave-one-coolant-inlet-out test:
for each chamber the operating-point and subcooling maps were refit on three of the four coolant
inlets and used to predict chip temperature at the withheld inlet, with the mechanistic chip model
held fixed. The 33-tube generalizes well, the held-out RMSE rising only from 3.8 to 5.0~K, and the
interior inlets (30 and 40~$^\circ$C) are predicted at in-sample quality; the edge inlets (20 and
50~$^\circ$C), which require extrapolating the linear map, are the hardest. The 42-tube degrades
more, from 4.9 to 9.0~K, consistent with its higher coolant gain (0.71 versus 0.26), which makes
the operating point depend strongly on coolant inlet, and with its poorly determined subcooling fit.

Second, a cross-chamber transfer of the plain-chip surface coefficient. Fit independently on the two
geometrically different chambers, $C_{sf}$ comes out at 0.0130 (33-tube) and 0.0121 (42-tube),
nearly identical and both close to the pooled value 0.0131. Applying the coefficient fit on one
chamber to the other raises the plain-chip RMSE by at most 0.2~K (33-tube 3.2 to 3.4~K; 42-tube 4.4
to 4.5~K). This directly supports treating $C_{sf}$ as a transferable surface--fluid property rather
than a per-chamber fit, which is the basis for predicting an untested chamber.

\begin{figure}[H]
\centering
\includegraphics[width=\textwidth]{mc_crossval.png}
\caption{Out-of-sample validation. Left: chip temperature predicted for a coolant inlet withheld
from the operating-point fit (held-out points only). Right: in-sample versus held-out error, with the
plain-chip $C_{sf}$ recovered independently on each chamber.}
\label{fig:crossval}
\end{figure}

\begin{table}[H]
\centering
\caption{Out-of-sample validation. Leave-one-coolant-inlet-out (LOCO) refits the operating-point and
subcooling maps on three coolant inlets and predicts the fourth; the transfer test applies a
plain-chip $C_{sf}$ fit on one chamber to the other.}
\label{tab:crossval}
\begin{tabular}{lcc}
\toprule
 & 33-tube & 42-tube\\
\midrule
In-sample chip-temperature RMSE (K) & 3.7 & 4.8\\
LOCO held-out chip-temperature RMSE (K) & 4.9 & 9.1\\
LOCO held-out $T_{sat}$ RMSE (K) & 7.4 & 15.3\\
\midrule
$C_{sf}$ fit on this chamber alone & 0.0130 & 0.0121\\
Plain-chip RMSE, native $C_{sf}$ (K) & 3.2 & 4.4\\
Plain-chip RMSE, transferred $C_{sf}$ (K) & 3.4 & 4.6\\
\bottomrule
\end{tabular}
\end{table}

\subsection{Origin of the 42-tube generalization gap}
\label{sec:gen42}
The leave-one-coolant-out gap between the chambers, 4.9~K held-out RMSE for the 33-tube against
9.1~K for the 42-tube, is worth isolating because both chambers share the same mechanistic chip
model; only the operating-point map differs. Figure~\ref{fig:gen42} traces the gap to its source,
and the chip physics is not it. Holding the true operating point fixed, the two chambers predict
surface temperature at the same quality; the difference lives entirely in how well the chamber
pressure can be predicted.

For the open 33-tube bundle the condenser-side resistance $(T_{liq}-T_{in})/Q$ is nearly
flux-independent, a median near 13~mK/W across the full flux range (Fig.~\ref{fig:gen42}a). For the
confined 42-tube bundle it is strongly flux-dependent, falling from roughly 150~mK/W at low flux to
about 45~mK/W near CHF, a slope of $-2.3$~mK/W per W/cm$^2$ against the 33-tube's $-0.07$. The
confined bundle is coolant-starved at low flux and progressively better swept as boiling intensifies,
so its pool-to-coolant resistance, and therefore the chamber pressure and saturation temperature it
sets, depend nonlinearly on heat load. The linear operating-point map $T_{sat}=a\,T_{in}+b\,Q+c$
cannot represent this: it fits the 42-tube saturation temperature with an 8.6~K RMS residual against
4.1~K for the 33-tube.

That operating-point error transmits almost one-to-one to surface temperature. Perturbing the
saturation temperature in the chip model gives $dT_{surf}/dT_{sat}\approx0.9$ for both chambers,
since in nucleate boiling the wall tracks saturation. Figure~\ref{fig:gen42}b confirms it: across the
held-out points the surface-temperature error follows the saturation-temperature error along the 0.9
line, and the 42-tube simply spans a far wider band of operating-point error, about $\pm35$~K against
$\pm15$~K. The three metrics line up accordingly (Fig.~\ref{fig:gen42}c): the 42-tube is worse by the
same factor of roughly two on the saturation-map residual (8.6 versus 4.1~K), the held-out saturation
error (13.6 versus 6.4~K), and the held-out surface-temperature RMSE (9.1 versus 4.9~K). The gap is
thus a limitation of representing a flux-dependent confined-bundle condenser with a load-linear map,
not of the boiling model. The implication for an untested confined geometry is concrete: the
operating-point sub-model should carry a mechanistic flux-dependent condenser resistance rather than a
linear-in-$Q$ fit, which is the single change most likely to bring 42-tube-class generalization in
line with the 33-tube.

\begin{figure}[H]
\centering
\includegraphics[width=\textwidth]{mc_gen42.png}
\caption{Origin of the 42-tube generalization gap. (a) Condenser-side resistance versus heat flux:
flat for the open 33-tube, strongly flux-dependent for the confined 42-tube (binned medians over the
data cloud). (b) Held-out surface-temperature error versus held-out saturation-temperature error;
both follow the $dT_{surf}/dT_{sat}\approx0.9$ line, the 42-tube over a wider range. (c) The 42-tube
is worse by about a factor of two on the saturation-map residual, the held-out saturation error, and
the held-out surface-temperature RMSE alike.}
\label{fig:gen42}
\end{figure}

\subsection{Comparison with baseline models}
\label{sec:baseline}
To show that the two main modeling choices earn their place, the full gray-box is compared against
two reduced alternatives (Table~\ref{tab:baseline}). The first replaces the per-surface boiling
coefficients with a single literature value $C_{sf}=0.013$ for both chips. On the plain chip this
is indistinguishable from the calibrated model (3.8 versus 3.8~K), confirming that the plain
calibration only recovers the handbook value; on the microchannel it nearly doubles the error
(8.8 versus 5.1~K), because the structured surface genuinely needs its own lower coefficient. The
second replaces the calibrated operating-point map with an ideal effectiveness--NTU condenser using
fully developed laminar coolant Nusselt number and the nominal tube area. Its effectiveness is only
$0.03$--$0.04$, so it over-predicts the saturation temperature by a median of $30~\mathrm{K}$ and
drives a few points above the critical point, which is nonphysical. Both comparisons reproduce the
paper's central claims: the microchannel requires a distinct surface coefficient, and the condenser
must be modeled as coolant-side limited rather than by ideal effectiveness--NTU.

\begin{table}[H]
\centering
\caption{Baseline comparison. Chip-temperature RMSE (K) for the full model versus a single-$C_{sf}$
evaporator; the effectiveness--NTU condenser is reported by its saturation-temperature error,
being nonphysical at high flux.}
\label{tab:baseline}
\begin{tabular}{lccc}
\toprule
Model & All & Plain & Microchannel\\
\midrule
Full gray-box & 4.4 & 3.8 & 5.1\\
Single $C_{sf}=0.013$, both chips & 6.4 & 3.8 & 8.8\\
\midrule
\multicolumn{4}{p{0.86\columnwidth}}{\footnotesize Effectiveness--NTU condenser: effectiveness
$0.03$--$0.04$; median $T_{sat}$ over-prediction $30~\mathrm{K}$; 3/261 points nonphysical
($>T_{crit}$).}\\
\bottomrule
\end{tabular}
\end{table}

\subsection{Condenser operating behavior}
\label{sec:condops}
Figure~\ref{fig:condops} shows how the condenser sets the operating point as the two control
parameters, coolant inlet temperature and chamber heat load, are varied. Removing the load
contribution with the fitted map isolates the coolant-inlet sensitivity (top left): saturation
temperature rises with coolant inlet at the coolant gain $a$, which is $0.26$ for the open 33-tube
and $0.71$ for the confined 42-tube, so the open chamber largely pins its saturation temperature
while the confined one tracks the coolant. Removing the inlet contribution isolates the load
sensitivity (top right), whose slope is the effective resistance $b$, again much steeper for the
42-tube. Pool subcooling falls with heat load (bottom left), and the condenser-side resistance
(bottom right) is roughly flux-independent at $\sim\!13~\mathrm{mK/W}$ for the 33-tube and
$\sim\!30$--$45~\mathrm{mK/W}$ for the 42-tube, consistent with the breakdown below.

\begin{figure}[H]
\centering
\includegraphics[width=\textwidth]{mc_condenser_ops.png}
\caption{Condenser operating behavior. Saturation temperature versus coolant inlet (coolant gain)
and heat load (effective resistance), pool subcooling versus load, and condenser-side resistance
versus flux. Markers are data pooled over both chips; lines are the operating-point model.}
\label{fig:condops}
\end{figure}

\subsection{Thermal resistance breakdown}
Figure~\ref{fig:resistance} gives the chip-side and condenser-side resistances versus heat flux to
CHF for each configuration, referenced to the bulk liquid temperature; values near CHF are in
Table~\ref{tab:resistance}. Referenced to saturation, the boiling heat transfer coefficient is
$\approx35~\mathrm{kW/m^2K}$ for the plain chip and up to $200~\mathrm{kW/m^2K}$ for the
microchannel.

\begin{figure}[H]
\centering
\includegraphics[width=0.49\textwidth]{res_33_plain.png}\hfill
\includegraphics[width=0.49\textwidth]{res_33_micro.png}\\[3pt]
\includegraphics[width=0.49\textwidth]{res_42_plain.png}\hfill
\includegraphics[width=0.49\textwidth]{res_42_micro.png}
\caption{Chip-side and condenser-side thermal resistance versus heat flux to CHF, referenced to
bulk liquid temperature. Markers are data (distinct symbol per side); lines are the model.}
\label{fig:resistance}
\end{figure}

\begin{table}[H]
\centering
\caption{Thermal resistance near CHF, referenced to bulk liquid, $R=\Delta T/Q$ in mK/W.}
\label{tab:resistance}
\begin{tabular}{lccc}
\toprule
Configuration & $R_{chip}$ & $R_{cond}$ & $R_{tot}$\\
\midrule
33-tube, plain        & 33 & 13 & 47\\
33-tube, microchannel & 11 & 13 & 24\\
42-tube, plain        & 36 & 46 & 81\\
42-tube, microchannel & 15 & 26 & 41\\
\bottomrule
\end{tabular}
\end{table}

The chip-side resistance is compared directly against experiment for both surfaces in
Figure~\ref{fig:chipres}. The microchannel's roughly threefold reduction is evident, from
$\sim\!38$ to $\sim\!13~\mathrm{mK/W}$ near CHF on the 33-tube, and the model tracks the
experimental chip-side resistance for both chips, the largest deviation being the 42-tube plain
configuration, consistent with its higher overall error.

\begin{figure}[H]
\centering
\includegraphics[width=\textwidth]{mc_chip_resistance.png}
\caption{Chip-side thermal resistance $R_{chip}=(T_{surf}-T_{bulk})/Q$ versus heat flux, model
versus experiment, referenced to bulk liquid. The microchannel cuts the chip-side resistance about
threefold relative to the plain chip; the model reproduces both surfaces.}
\label{fig:chipres}
\end{figure}

\section{Discussion}
The central physical finding is that the condenser is limited by coolant-side
developing-laminar convection rather than by condensation, and that its conductance improves
with operating temperature. This produces a saturation-temperature gain on coolant inlet below
unity and explains why the open 33-tube self-regulates while the confined 42-tube, with smaller,
denser, more submerged tubes, climbs. The resistance split referenced to bulk liquid shows that
the chip side, which carries the pool subcooling as well as the wall superheat, is the larger
resistance for the open 33-tube, whereas the condenser side dominates for the confined 42-tube.
The microchannel cuts the chip-side resistance roughly threefold, after which the condenser side
is comparable (33-tube) or dominant (42-tube). The design consequence is explicit: surface
enhancement of the chip is effective, but for the resistive 42-tube the condenser is the
bottleneck, and further gains require improving it rather than the chip. The model also separates
transferable quantities ($C_{sf}$, $A_r$) from chamber-specific ones ($F_{nc}$ and the
operating-point coefficients), the latter being physically interpretable.

\subsection{Predictive scope for untested geometry}
\label{sec:scope}
The two sides of the model are predictive to different degrees, which determines what can be
forecast for hardware that has not been tested.

\textit{New chip or area geometry: predictive.} The evaporator is fully mechanistic. The
surface coefficient $C_{sf}$ is a transferable surface--fluid property (Sections~\ref{sec:coeff},
\ref{sec:crossval}): it transfers across chambers at under $0.2~\mathrm{K}$ and matches handbook
values by surface finish. The area augmentation $A_r$ is computed from the channel geometry, not
fitted. A new boiling surface therefore needs only its $C_{sf}$, assigned from its finish or pinned
by a single boiling curve, together with its geometric $A_r$; the full $T_{surf}(q'')$ curve then
follows with no chamber-specific evaporator fitting. The exception is the burnout limit: the
microchannel CHF is an additive \emph{measured} increment and the mechanistic CHF model is an
uncalibrated placeholder, so the curve is predictable up to CHF but the CHF value itself needs a
measurement or a dedicated closure.

\textit{New condenser geometry: ranking predictable, absolute level needs one calibration.}
The coolant-side resistance follows from geometry through the developing-laminar internal-flow
correlation, and it sets the right order of magnitude and ordering
(Table~\ref{tab:condmech}, Figure~\ref{fig:condscope}): the measured condenser-side resistance is
$13.5~\mathrm{mK/W}$ for the 33-tube against a geometric estimate of $17.3$, and $42~\mathrm{mK/W}$
for the 42-tube against $25.7$. The two chambers bracket the ideal, the open 33-tube does
slightly better than full-area developing-laminar convection (additional shell-side mechanisms),
the confined 42-tube does worse (only part of the bundle actively condenses), so the ratio of
real to ideal resistance is a layout-dependent factor ($0.8$ open, $1.6$ confined) that tube
dimensions alone do not pin. As a result the absolute operating point (the $T_{sat}$ and subcooling
maps) is currently fit per condenser from one dataset; once fit, the model predicts across coolant
inlets, fluxes, and chips for that condenser. Making the condenser side fully geometry-predictive, replacing the operating-point map with an active-area closure that returns $(a,b,c)$ from
geometry and fill, is the main open step, and it requires at least a third characterized
condenser to fit and test.

\textit{The 66-tube.} Its coolant-side resistance is the lowest of the three by geometry
($\sim\!11~\mathrm{mK/W}$ at the 33-tube coolant flow, owing to its large area and short tubes),
so on condenser hardware alone it should hold the lowest $T_{sat}$ for a given load and run the
coolest chip. Being open like the 33-tube, its real-to-ideal factor is expected near the
33-tube's, which bounds its operating point close to the geometric estimate. This remains a
prediction until one 66-tube dataset is taken, which is the recommended next measurement.

\begin{table}[H]
\centering
\caption{Condenser-side thermal resistance: geometric developing-laminar coolant-side estimate
versus measured. The real-to-ideal ratio is layout-dependent, so it is not fixed by tube
dimensions alone.}
\label{tab:condmech}
\begin{tabular}{lcccc}
\toprule
Condenser & $Re$ & Geometric $R$ (mK/W) & Measured $R$ (mK/W) & real/ideal\\
\midrule
33-tube (open)     & 1200 & 17.3 & 13.5 & 0.78\\
42-tube (confined) & 985  & 25.7 & 42.0 & 1.63\\
66-tube (open)     & 1050 & 11.1 & $-$ & $-$\\
\bottomrule
\end{tabular}
\end{table}

\begin{figure}[H]
\centering
\includegraphics[width=0.78\textwidth]{mc_condscope.png}
\caption{Condenser-side resistance from geometry versus measurement. The developing-laminar
coolant-side estimate sets the ranking and order of magnitude; a layout-dependent factor separates
ideal from real. The 66-tube has the lowest geometric resistance but is not yet measured.}
\label{fig:condscope}
\end{figure}

\subsection{Predicted 66-tube performance}
\label{sec:pred66}
The 66-tube operating point is not measured, but its one geometric difference, a lower
condenser-side resistance, can be propagated through a construction that is first validated on the
33-tube. The pool temperature is set by the condenser, $T_{bulk}=T_{in}+Q\,R_{cond}$, with $R_{cond}$
the developing-laminar coolant-side estimate scaled by the open real-to-ideal factor; the saturation
temperature is $T_{sat}=T_{bulk}+\Delta T_{sub}$ with subcooling from the open-bundle map; and the
chip temperature follows from the transferable boiling model. Applied to the 33-tube this reproduces
the calibrated model to within $0.2~\mathrm{K}$, so the same construction is used for the 66-tube
with its geometric resistance of $8.7~\mathrm{mK/W}$ ($11.1\times0.78$).

Tables~\ref{tab:geom66} and \ref{tab:pred66} give the geometry and the predicted operating point at
CHF. Rejecting heat through roughly half the condenser-side resistance of the 33-tube, the 66-tube
holds a cooler pool and a lower saturation temperature for the same load: the chip runs about
$4~\mathrm{K}$ cooler on the plain surface and $8~\mathrm{K}$ cooler on the microchannel at CHF,
the $\sim\!6~\mathrm{K}$ drop in saturation temperature outweighing a small rise in wall superheat
at the lower pressure. The condenser-side resistance falls from $0.0135$ to $0.0087~\mathrm{K/W}$
while the chip-side resistance, fixed by the same surface and flux, is nearly unchanged. These are
predictions under stated assumptions, CHF and pool subcooling taken from the open 33-tube, that
one 66-tube dataset would replace and test directly.

\begin{table}[H]
\centering
\caption{Geometry of the three condensers. Internal and external areas are $\pi\,(ID)\,L\,N$ and
$\pi\,(OD)\,L\,N$. The 66-tube coolant flow is assumed equal to the 33-tube loop ($^\ast$).}
\label{tab:geom66}
\begin{tabular}{lccc}
\toprule
 & 33-tube & 42-tube & 66-tube\\
\midrule
Tubes $N$ & 33 & 42 & 66\\
Outer diameter (mm) & 4.76 & 3.18 & 2.00\\
Inner diameter (mm) & 3.14 & 1.39 & 1.80\\
Tube length (mm) & 95 & 62.6 & 95\\
Internal area (cm$^2$) & 309 & 115 & 355\\
External area (cm$^2$) & 469 & 262 & 394\\
Coolant flow (g/s) & 78 & 36 & 78$^\ast$\\
Coolant $Re$ & 1200 & 985 & 1050\\
Manifolds & open ext. & in-chamber & open ext.\\
Geometric coolant-side $R$ (mK/W) & 17.3 & 25.7 & 11.1\\
\bottomrule
\end{tabular}
\end{table}

\begin{table}[H]
\centering
\caption{Predicted 66-tube performance at CHF (coolant inlet $30~^\circ$C), with the 33-tube for
comparison. CHF is taken equal to the 33-tube; the resistance split is referenced to bulk liquid.}
\label{tab:pred66}
\begin{tabular}{lcccc}
\toprule
 & \multicolumn{2}{c}{66-tube (predicted)} & \multicolumn{2}{c}{33-tube}\\
 & plain & microch. & plain & microch.\\
\midrule
$q''_{CHF}$ (W/cm$^2$)        & 114  & 174  & 114  & 174\\
Heat load $Q$ (W)            & 1259 & 1921 & 1259 & 1921\\
$T_{sat}$ ($^\circ$C)        & 53.9 & 55.2 & 60.0 & 64.5\\
Subcooling (K)               & 13.0 & 8.5  & 13.0 & 8.5\\
$T_{surf}$ at CHF ($^\circ$C)& 85.6 & 69.8 & 89.4 & 77.5\\
\midrule
$R_{chip}$ (K/W) & 0.0355 & 0.0121 & 0.0337 & 0.0112\\
$R_{cond}$ (K/W) & 0.0087 & 0.0087 & 0.0135 & 0.0135\\
$R_{tot}$ (K/W)  & 0.0442 & 0.0207 & 0.0472 & 0.0247\\
\bottomrule
\end{tabular}
\end{table}

Figure~\ref{fig:bc66} shows the predicted boiling curve obtained by sweeping the same construction
to CHF. The left panel is the boiling curve proper, heat flux against wall superheat: below boiling
onset the surface sits in subcooled single-phase convection, and above it the nucleate branch rises
steeply, with the microchannel reaching $174~\mathrm{W/cm^2}$ at about $15~\mathrm{K}$ superheat
against $32~\mathrm{K}$ for the plain chip. The right panel plots flux against chip temperature with
the 33-tube dashed for reference, showing the 66-tube holding a lower chip temperature at every flux
because of its lower condenser-side resistance.

\begin{figure}[H]
\centering
\includegraphics[width=\textwidth]{mc_boilingcurve_66.png}
\caption{Predicted 66-tube boiling curve at $30~^\circ$C coolant inlet. Left: heat flux versus wall
superheat (single-phase below onset, nucleate above) to CHF. Right: heat flux versus chip surface
temperature, with the 33-tube (dashed) for comparison. Curves use the geometry-based construction
validated against the 33-tube; CHF is taken from the 33-tube pending a 66-tube measurement.}
\label{fig:bc66}
\end{figure}

\section{Conclusions}
\begin{enumerate}
\item A reduced-order model predicts chip surface temperature from heat flux, coolant inlet
temperature, chip surface, and condenser geometry, with no measured operating conditions, to a
root-mean-square error of 4.4~K (median 4\% of absolute temperature; 5\% on the rise above coolant
in the high-flux regime).
\item The microchannel chip is captured by the same boiling machinery with a fixed area
augmentation ($A_r=2.17$) and a halved, transferable surface coefficient ($C_{sf}=0.0067$ versus
$0.0131$); its CHF is the plain CHF plus a constant $\sim\!60~\mathrm{W/cm^2}$.
\item The condenser is coolant-side limited, and its coolant gain ($0.26$ open, $0.71$ confined)
is a fingerprint of the chamber.
\item Referenced to bulk liquid, the condenser side dominates the chip-to-coolant resistance for
the confined 42-tube and is comparable to the chip side for the 33-tube once the microchannel is
used, identifying the condenser as the target for further improvement.
\item Out-of-sample tests support the predictive claim: the plain-chip surface coefficient fit
independently on the two chambers agrees (0.0130 versus 0.0121) and transfers between them at a cost
below 0.2~K, while a leave-one-coolant-inlet-out test predicts chip temperature at a withheld inlet
with 5.0~K RMSE on the 33-tube (9.0~K on the 42-tube, limited by its stronger coolant-inlet
dependence).
\item The model is predictive on the evaporator side, a new chip surface is forecast from a
transferable $C_{sf}$ and a geometric area augmentation, up to a CHF limit that still needs
measurement, whereas the condenser side predicts the resistance ranking from geometry but
requires one calibration dataset to fix the absolute operating point, which sets the agenda for a
geometry-based condenser closure.
\end{enumerate}

\section{Limitations and future work}
The operating-point coefficients are a calibrated per-chamber map, not a first-principles closure,
partly because the coolant flow carries measurement uncertainty and the level dynamics that grow
the condensing area are not resolved from geometry alone. Subcooling is well predicted on the
33-tube and essentially scatter on the 42-tube. CHF is anchored to measured burnout. The immediate
extension is a single run on the 66-tube condenser to fix its operating-point coefficients. The
HFE-7000 path in the model is at present a literature-anchored prediction, with smooth-copper boiling
coefficients rather than a chamber-data fit, so extending the model to a validated second working
fluid, together with a mechanistic CHF closure, are longer-term goals that require an HFE dataset.

\section{Next actions}
The following are ordered by how much each converts a standing assumption into validated physics.
\begin{enumerate}
\item \textbf{Take one 66-tube dataset.} A single sweep on the existing 66-tube condenser validates
the prediction in Section~\ref{sec:pred66}, fixes its operating-point coefficients, and provides the
third characterized condenser needed to begin a geometry-based condenser closure. This is the single
highest-value measurement.
\item \textbf{Resolve whether CHF is burnout or a surface-temperature ceiling.} The same plain chip
reaches 114~W/cm$^2$ on the 33-tube and 65 on the 42-tube at comparable subcooling and similar
surface temperature, which suggests a system limit rather than a local boiling crisis. Settling this
from the existing records determines whether the 66-tube and HFE CHF values should be inherited,
raised (lower-resistance condenser reaches a ceiling at higher flux), or modeled mechanistically.
\item \textbf{Bring in HFE-7000.} When data is available, refit $C_{sf}$ for HFE on copper (expected
near scored-copper values, with a much lower CHF), recompute the boiling curve, and test whether the
coolant-side-limited operating-point structure holds for a second fluid. The El-Genk correlations
give an independent check on the boiling coefficient and the subcooling dependence.
\item \textbf{Complete the measurement-uncertainty section.} Insert the instrument specifications
(thermocouple type and accuracy, sensor-to-surface distance, heat-flux and pressure uncertainty),
then add error bars to the data figures so the model error can be read against the noise floor.
\item \textbf{Document the physical apparatus.} Add a labeled apparatus schematic, the full
experimental test matrix, and a thermophysical property table for the working fluids.
\item \textbf{Develop the geometry-based condenser closure.} Replace the empirical operating-point
map with an active-area model that returns the coefficients from condenser geometry and liquid fill,
which is the step that makes the condenser side fully predictive. This requires the 66-tube data and,
ideally, the fill-ratio sweeps.
\item \textbf{Strengthen the 42-tube subcooling closure}, currently essentially scatter, and add a
mechanistic CHF model to remove the last measured anchor.
\end{enumerate}



\clearpage
\part{Three-Dimensional VOF Phase-Change Framework}

\section{Overview and scope of the CFD framework}
Two-phase immersion cooling is a leading candidate for the rising thermal loads of
data-center processors. The apparatus modeled here is a closed chamber in which a
heated chip boils a subcooled pool and the vapor is condensed on a tube bundle that is
partially submerged in the same pool, returning condensate without an external
condenser loop~\cite{of-shukla2024}. The platform has been characterized across two
evaporator surfaces (plain copper and open microchannel), two condenser bundles (33
tube and 42 tube), and a range of coolant inlet temperatures; a validated reduced-order
model predicts the chip surface temperature from heat flux for these configurations.
The present work develops the complementary high-fidelity route: a three-dimensional
VOF simulation with phase change that resolves the boiling and condensation fields on
the real geometry.

This document specifies that framework completely and reproducibly, and it assembles
the experimental reference, coolant operating conditions, measured boiling curves, and
thermal resistances for both condensers, that the simulation is required to reproduce.
Following the integrity standard of the broader project, we keep a hard line between
three sources: measured data, the validated reduced-order model, and the CFD. The
boiling-curve and resistance comparisons in Section~\ref{of-sec:reference} are
\emph{experiment versus the reduced-order model}; they are the CFD's target, not CFD
output. The CFD field contours and the CFD-to-experiment overlay require the converged
run of Section~\ref{of-sec:post} and are not fabricated here.

\section{Apparatus geometry}
\label{of-sec:geom}

\subsection{Chamber and computational domain}
The domain is the chamber interior, $80 \times 80 \times \SI{42}{mm}$ ($x$~width,
$y$~tube axis, $z$~height), gravity along $-z$ (Fig.~\ref{of-fig:domain}). The 42 condenser
tubes are imported as a triangulated surface at their measured CAD positions; the bundle
straddles the liquid level so the lower rows are submerged and the upper rows sit in the
vapor space. The heated chip ($34.5\times\SI{32}{mm}$, \SI{11.04}{cm^2} footprint) is
centered on the floor, and the pool fills the lower \SI{16.8}{mm} (about 40\%).
Figure~\ref{of-fig:geom} gives the dimensioned cross-section with the real staggered tube
layout.

\begin{figure}[H]
\centering
\begin{subfigure}{0.56\textwidth}\includegraphics[width=\textwidth]{fig_domain.png}
\caption{Isometric.}\end{subfigure}\hfill
\begin{subfigure}{0.42\textwidth}\includegraphics[width=\textwidth]{fig_domain_front.png}
\caption{Front ($x$--$z$).}\end{subfigure}
\caption{Computational domain: chamber (outline), 42-tube bundle (grey), heated chip
(red), subcooled pool (blue) at 40\% fill. The bundle straddles the interface.}
\label{of-fig:domain}
\end{figure}

\begin{figure}[H]
\centering\includegraphics[width=0.80\textwidth]{fig_geom_dims.png}
\caption{Dimensioned $x$--$z$ cross-section with the measured 42-tube layout; the dashed
line is the initial interface at \SI{16.8}{mm}.}
\label{of-fig:geom}
\end{figure}

\subsection{Evaporator surfaces}
Both chips share the \SI{11.04}{cm^2} copper footprint. The plain chip is a smooth
boiling surface; the enhanced chip carries open microchannels of the Cooke--Kandlikar
type, \SI{381}{\micro m}-wide channels, \SI{400}{\micro m} deep, separated by
\SI{250}{\micro m} fins (Fig.~\ref{of-fig:micro}), which raise the wetted-to-footprint area
by a factor $A_r=2.17$ and provide rewetting pathways. In the CFD the plain surface is a
flat heated wall; resolving or sub-grid-modeling the microchannel surface is a planned
extension (Section~\ref{of-sec:next}).

\begin{figure}[H]
\centering
\begin{subfigure}{0.50\textwidth}\includegraphics[width=\textwidth]{fig_micro_3d.png}
\caption{Surface portion (6\,mm).}\end{subfigure}\hfill
\begin{subfigure}{0.48\textwidth}\includegraphics[width=\textwidth]{fig_micro_section.png}
\caption{Cross-section with dimensions.}\end{subfigure}
\caption{Open-microchannel evaporator surface (Cooke--Kandlikar geometry).}
\label{of-fig:micro}
\end{figure}

\subsection{Condenser bundles}
Two bundles are studied (Fig.~\ref{of-fig:cond}, Table~\ref{of-tab:cond}). The 33-tube bundle
uses fewer, larger tubes (OD \SI{4.76}{mm}) with open manifolds; the 42-tube bundle uses
more, smaller tubes (OD \SI{3.175}{mm}) in a confined in-chamber manifold. The CFD case
documented here is the 42-tube geometry; a 33-tube case is a parallel build
(Section~\ref{of-sec:next}).

\begin{figure}[H]
\centering\includegraphics[width=0.92\textwidth]{fig_condensers.png}
\caption{The two condenser bundles to scale: 33-tube (OD \SI{4.76}{mm}; layout
schematic, CAD positions not available) and 42-tube (OD \SI{3.175}{mm}, from CAD).}
\label{of-fig:cond}
\end{figure}

\begin{table}[H]
\centering\caption{Condenser bundle geometry.}\label{of-tab:cond}
\begin{tabular}{lccccc}
\toprule
Bundle & Tubes & OD (mm) & ID (mm) & $L$ (mm) & Layout\\
\midrule
33-tube & 33 & 4.76 & 3.14 & 95.0 & open manifold\\
42-tube & 42 & 3.175 & 1.39 & 62.6 & confined\\
\bottomrule
\end{tabular}
\end{table}

\section{Governing equations}
\label{of-sec:equations}
The solver \texttt{interCondensatingEvaporatingFoam} advances a single incompressible
two-phase (water/vapor) mixture with volume-of-fluid interface capturing, buoyancy,
surface tension, and a volumetric phase-change source. The flow is treated as laminar
(justified in Section~\ref{of-sec:dimensionless}). A single momentum and energy field is
solved for the mixture, with the two phases distinguished by the liquid volume fraction
$\alpha\equiv\texttt{alpha.water}\in[0,1]$.

\subsection{Two-phase VOF transport}
Mass conservation of each phase, with interconversion $\dot m$ (mass per unit volume per
unit time, positive for evaporation), gives the indicator-function transport equation
solved with an interface-compression flux $\mathbf U_r$ that counteracts numerical
smearing,
\begin{equation}
\frac{\partial \alpha}{\partial t}+\nabla\!\cdot(\mathbf{U}\alpha)
+\nabla\!\cdot\!\big(\mathbf{U}_r\,\alpha(1-\alpha)\big)
=\frac{\dot m_c-\dot m_e}{\rho_l},
\qquad
\mathbf U_r=c_\alpha\,|\mathbf U|\,\frac{\nabla\alpha}{|\nabla\alpha|},
\label{of-eq:alpha}
\end{equation}
where $c_\alpha=1$ sets the compression strength. The compression term acts only in the
interface band ($0<\alpha<1$) and vanishes in the bulk. Because mass is exchanged between
phases of unequal density, the mixture velocity is not divergence-free; continuity becomes
\begin{equation}
\nabla\!\cdot\mathbf U=\dot m\!\left(\frac{1}{\rho_v}-\frac{1}{\rho_l}\right),
\label{of-eq:cont}
\end{equation}
so evaporation is a local volume source and condensation a local sink. Mixture properties
are volume-fraction weighted,
\begin{equation}
\rho=\alpha\rho_l+(1-\alpha)\rho_v,\quad
\mu=\alpha\mu_l+(1-\alpha)\mu_v,\quad
(\rho c_p)=\alpha\rho_l c_{p,l}+(1-\alpha)\rho_v c_{p,v}.
\end{equation}

\subsection{Momentum and surface tension}
Momentum is solved for the mixture in the buoyant $p_{rgh}$ form, where
$p_{rgh}=p-\rho\,\mathbf g\!\cdot\!\mathbf x$ removes the hydrostatic head from the
pressure solution and improves conditioning,
\begin{equation}
\frac{\partial(\rho\mathbf U)}{\partial t}+\nabla\!\cdot(\rho\mathbf U\mathbf U)
=-\nabla p_{rgh}-(\mathbf g\!\cdot\!\mathbf x)\nabla\rho
+\nabla\!\cdot\!\big[\mu(\nabla\mathbf U+\nabla\mathbf U^{\mathsf T})\big]
+\mathbf F_{st}.
\label{of-eq:mom}
\end{equation}
The second term is the buoyancy force carried by density gradients (the engine of natural
convection in the subcooled pool). Surface tension enters as a continuum surface force
\cite{of-brackbill1992} localized at the interface,
\begin{equation}
\mathbf F_{st}=\sigma\kappa\nabla\alpha,\qquad
\kappa=-\nabla\!\cdot\!\left(\frac{\nabla\alpha}{|\nabla\alpha|}\right),
\end{equation}
with $\kappa$ the interface curvature. Velocity is coupled to pressure through
Eq.~\eqref{of-eq:cont} in a PISO loop (Section~\ref{of-sec:numerics}).

\subsection{Energy and the latent coupling}
Energy is solved in temperature form for the mixture, with conduction through an effective
conductivity and a volumetric latent sink/source at the interface,
\begin{equation}
\frac{\partial(\rho c_p T)}{\partial t}+\nabla\!\cdot(\rho c_p\mathbf U T)
=\nabla\!\cdot(k_{\mathrm{eff}}\nabla T)-\dot m\,h_{fg},
\qquad
k_{\mathrm{eff}}=\alpha k_l+(1-\alpha)k_v.
\label{of-eq:energy}
\end{equation}
Evaporation removes $\dot m_e h_{fg}$ from the liquid at the heated surface; condensation
releases $\dot m_c h_{fg}$ at the cold tubes. The same effective conductivity defines the
chip heat-flux probe used to build the simulated boiling curve (Section~\ref{of-sec:post}).

\subsection{Phase-change closure}
The interfacial mass flux is closed with the constant-coefficient (Lee) model
\cite{of-lee1980}, in which the rate is proportional to the local departure from saturation,
driving the metastable cells back toward $T_{sat}$ (Fig.~\ref{of-fig:pc}),
\begin{equation}
\dot m_e=C_E\,\alpha\,\rho_l\,\frac{T-T_{sat}}{T_{sat}}\;\;(T>T_{sat}),
\qquad
\dot m_c=C_C\,(1-\alpha)\,\rho_v\,\frac{T_{sat}-T}{T_{sat}}\;\;(T<T_{sat}),
\label{of-eq:lee}
\end{equation}
where $C_E$ (\texttt{coeffE}) and $C_C$ (\texttt{coeffC}) are relaxation frequencies
(\si{\per\second}). The coefficient sets an interfacial relaxation time
$\tau\sim1/(C\,\Delta T/T_{sat})$: too small and the interface superheats unphysically;
too large and the equations stiffen and the time step collapses. There is no universal
value, $C_E$ and $C_C$ are mesh- and resolution-dependent and are calibrated against the
measured boiling curve (Section~\ref{of-sec:reference}); this is the central empirical content
of the model and the reason a calibration campaign, not a single run, is required.

\begin{figure}[H]
\centering\includegraphics[width=0.84\textwidth]{fig_phasechange.png}
\caption{Phase-change closure: evaporation where the cell is superheated ($T>T_{sat}$) and
condensation where it is subcooled ($T<T_{sat}$), each at a rate proportional to the
departure from saturation. The interface is the $\alpha=0.5$ surface.}
\label{of-fig:pc}
\end{figure}

\subsection{Flow regime and dimensionless groups}
\label{of-sec:dimensionless}
Table~\ref{of-tab:dim} collects the governing dimensionless groups for water at the operating
point. Three features shape the modeling choices. First, the pool Rayleigh number
($\mathrm{Ra}\sim9\times10^6$, based on the liquid depth and the chip-to-pool difference)
sits well below the $\sim10^9$ transition for a heated horizontal layer, so the bulk
buoyant flow is laminar; the vigorous agitation in the chamber is driven by boiling and
bubble departure rather than by turbulent natural convection, which supports the laminar
mixture treatment with the interface resolved. Second, the Bond and confinement numbers
based on tube diameter ($\mathrm{Bo}=1.5$--$3.4$, $\mathrm{Co}=L_c/D=0.5$--$0.8$) show the
capillary length $L_c=\sqrt{\sigma/\rho_l g}\approx\SI{2.6}{mm}$ is comparable to the tube
spacing, so capillarity and confinement matter at the bundle, the smaller-bore confined
42-tube bundle is more strongly affected than the open 33-tube. Third, the liquid-to-vapor
density ratio is large ($\rho_l/\rho_v\approx3300$) and the superheat Jakob number is small
($\mathrm{Ja}=c_{p,l}\Delta T_{sup}/h_{fg}\approx0.04$); the large density ratio is the main
source of numerical stiffness (the volume source in Eq.~\eqref{of-eq:cont} and the property
jump), which is why the time step is interface-Courant limited. The chamber operates in
nucleate boiling up to CHF (Fig.~\ref{of-fig:regime}); the model is not intended for the
transition or film-boiling branches.

\begin{table}[H]
\centering\caption{Dimensionless groups (water, operating point: $\Delta T_{sup}\approx
\SI{20}{K}$, $\Delta T_{sub}\approx\SI{15}{K}$, liquid depth \SI{16.8}{mm}).}
\label{of-tab:dim}
\small
\begin{tabular}{llcl}
\toprule
Group & Definition & Value & Role\\
\midrule
Capillary length & $L_c=\sqrt{\sigma/\rho_l g}$ & \SI{2.6}{mm} & bubble / meniscus scale\\
Prandtl (liquid) & $\mu c_p/k$ & 2.7 & momentum vs thermal diffusion\\
Jakob (superheat) & $c_{p,l}\Delta T_{sup}/h_{fg}$ & 0.036 & sensible vs latent at the chip\\
Jakob (subcooling) & $c_{p,l}\Delta T_{sub}/h_{fg}$ & 0.027 & condensation driving\\
Density-ratio Jakob & $\rho_l c_{p,l}\Delta T_{sup}/(\rho_v h_{fg})$ & 117 & vapour generation intensity\\
Bond (42 / 33-tube) & $\rho_l g\,D^2/\sigma$ & 1.5 / 3.4 & gravity vs capillarity at a tube\\
Confinement (42 / 33) & $L_c/D$ & 0.82 / 0.55 & bundle confinement\\
Rayleigh (pool) & $g\beta\Delta T H^3/\nu\alpha_{th}$ & $9\times10^6$ & laminar bulk convection\\
Density ratio & $\rho_l/\rho_v$ & 3300 & numerical stiffness\\
\bottomrule
\end{tabular}
\end{table}

\begin{figure}[H]
\centering\includegraphics[width=0.72\textwidth]{fig_regime.png}
\caption{Pool-boiling regime map (schematic). The chamber operates on the nucleate-boiling
branch up to the measured CHF; the constant-coefficient closure is calibrated for this
branch, not for transition or film boiling.}
\label{of-fig:regime}
\end{figure}

\section{Mesh generation}
\label{of-sec:mesh}
A structured background grid (\texttt{blockMesh}, $40\times40\times21$, \num{33600}
cells) is refined with \texttt{snappyHexMesh}: surface refinement and feature snapping on
the tubes, refinement boxes over the chip and bundle, and prism layers. The chip patch is
cut from the floor. Body-fitted meshing of all 42 tubes is verified
(Fig.~\ref{of-fig:meshref}): snapping yields a complete grid of \num{299656} cells passing
\texttt{checkMesh} (max skewness \num{1.57}, 100\% layer addition), confirming the
limitation to a fully refined grid is cost, not mesh validity (Table~\ref{of-tab:mesh}).

\begin{figure}[H]
\centering
\begin{subfigure}{0.99\textwidth}\includegraphics[width=\textwidth]{refined_mesh_slice.png}
\caption{Mid-plane slice: body-fitted cells with graded refinement and the chip box.}
\end{subfigure}\\[4pt]
\begin{subfigure}{0.78\textwidth}\includegraphics[width=\textwidth]{refined_mesh_zoom.png}
\caption{Detail at the tube walls.}\end{subfigure}
\caption{Verified body-fitted snapped mesh (\num{299656} cells, \texttt{checkMesh} OK).}
\label{of-fig:meshref}
\end{figure}

\begin{table}[H]
\centering\caption{Mesh states.}\label{of-tab:mesh}
\begin{tabular}{lccl}
\toprule
Mesh & Construction & Cells & Notes\\
\midrule
Background & \texttt{blockMesh} & \num{33600} & $\sim$\SI{2}{mm}\\
Demonstration & castellation-only & \num{74044} & staircased, robust\\
Verified snapped & snap, level $(1\,2)$ & \num{299656} & \texttt{checkMesh} OK, skew \num{1.57}\\
Production target & shipped dict & $\sim$few$\times10^6$ & \SIrange{0.1}{0.3}{mm}, snap + layers (HPC)\\
\bottomrule
\end{tabular}
\end{table}

\section{Boundary, initial, and coolant operating conditions}
\label{of-sec:bc}
The chip is a fixed-temperature wall and the condenser tubes a colder fixed-temperature
wall; remaining walls are adiabatic no-slip, $T_{sat}=\SI{338}{K}$, and the pool is
initialized below \SI{16.8}{mm} (Fig.~\ref{of-fig:bc}, Table~\ref{of-tab:bc}). These wall
temperatures are representative; the calibration matches the simulated boiling curve to
the measured one (Section~\ref{of-sec:reference}).

\begin{figure}[H]
\centering\includegraphics[width=0.72\textwidth]{fig_bc.png}
\caption{Boundary and initial conditions on the $x$--$z$ cross-section.}
\label{of-fig:bc}
\end{figure}

\begin{table}[H]
\centering\caption{Field boundary conditions.}\label{of-tab:bc}
\begin{tabular}{llll}
\toprule
Patch & $T$ & $\mathbf U$ & $p_{rgh}$\\
\midrule
chip & fixedValue \SI{358}{K} & noSlip & fixedFluxPressure\\
condenser & fixedValue \SI{305}{K} & noSlip & fixedFluxPressure\\
floor / walls / top & zeroGradient & noSlip & fixedFluxPressure\\
\midrule
\multicolumn{4}{l}{Internal $T=\SI{325}{K}$, $\mathbf U=0$, $\alpha$ from
\texttt{setFields}; $T_{sat}=\SI{338}{K}$.}\\
\bottomrule
\end{tabular}
\end{table}

\noindent\textbf{Condenser coolant operating point (measured).} The condenser tubes
carry single-phase water coolant; its measured operating conditions set the cold-side
boundary that the fixed-temperature wall idealizes. Table~\ref{of-tab:coolant} reports the
inlet and outlet temperatures, mass flow, the resulting per-tube velocity and Reynolds
number, and the chamber pressure for both condensers across the four coolant set points.
The flow is laminar and developing in both bundles ($\mathrm{Re}\sim700$--$1600$); the
42-tube bundle runs at higher velocity (smaller bore) and lower chamber pressure at the
warm set points. The coolant temperature rise $T_{out}-T_{in}$ is the calorimetric basis
for the condenser duty $Q=\dot m c_p (T_{out}-T_{in})$.

\begin{table}[H]
\centering\caption{Measured condenser coolant conditions (per coolant set point; values
are run-averaged). $V$ and $\mathrm{Re}$ are per-tube, evaluated at $T_{in}$.}
\label{of-tab:coolant}
\footnotesize
\begin{tabular}{llcccccc}
\toprule
Bundle, chip & Set & $T_{in}$ (\si{\celsius}) & $T_{out}$ (\si{\celsius}) &
$\Delta T$ (K) & $\dot m$ (g/s) & $V$ (m/s) & $\mathrm{Re}$\\
\midrule
33, plain & 20 & 20.0 & 21.7 & 1.7 & 60.0 & 0.235 & 735\\
          & 30 & 29.2 & 31.0 & 1.8 & 60.0 & 0.236 & 910\\
          & 40 & 39.4 & 41.0 & 1.6 & 60.0 & 0.237 & 1116\\
          & 50 & 48.3 & 49.9 & 1.6 & 60.0 & 0.237 & 1310\\
33, micro & 20 & 20.2 & 23.3 & 3.1 & 62.5 & 0.245 & 771\\
          & 40 & 40.9 & 44.4 & 3.4 & 62.5 & 0.247 & 1197\\
42, plain & 20 & 20.0 & 22.1 & 2.1 & 38.8 & 0.610 & 845\\
          & 30 & 30.0 & 32.5 & 2.5 & 38.8 & 0.611 & 1061\\
          & 40 & 40.1 & 42.2 & 2.2 & 38.8 & 0.614 & 1298\\
          & 50 & 50.0 & 52.0 & 2.0 & 38.8 & 0.616 & 1549\\
42, micro & 20 & 20.0 & 23.5 & 3.5 & 41.0 & 0.644 & 893\\
          & 30 & 29.9 & 34.5 & 4.5 & 38.8 & 0.611 & 1060\\
          & 40 & 40.0 & 44.3 & 4.3 & 38.8 & 0.614 & 1296\\
          & 50 & 49.9 & 54.9 & 5.0 & 38.8 & 0.616 & 1547\\
\bottomrule
\end{tabular}\\[2pt]
\footnotesize Chamber pressure ranges \SIrange{72}{85}{kPa} (33-tube) and
\SIrange{37}{85}{kPa} (42-tube), falling at the warm set points as $T_{sat}$ drops.
\end{table}

\section{Material properties and numerics}
\label{of-sec:numerics}
Properties are in Table~\ref{of-tab:props}; the phase-change coefficients are the
uncalibrated placeholders $C_C=1.0$, $C_E=0.1$. Time integration is Euler with an
adjustable step ($\mathrm{Co}\le0.4$); convection uses \texttt{linearUpwind} (momentum,
energy) and \texttt{vanLeer} with interface compression for $\alpha$; pressure uses
PCG/DIC, $\alpha$ uses MULES, within a PISO loop (Table~\ref{of-tab:run}).

\begin{table}[H]
\centering
\begin{minipage}[t]{0.54\textwidth}\centering\caption{Material properties.}\label{of-tab:props}\small
\begin{tabular}{lccc}
\toprule
Property & Water & Vapor & Unit\\\midrule
$\rho$ & 980 & 0.30 & \si{kg/m^3}\\
$\nu$ & \num{4.4e-7} & \num{4.0e-4} & \si{m^2/s}\\
$c_p$ & 4190 & 2030 & \si{J/kg/K}\\
$\kappa$ & 0.66 & 0.025 & \si{W/m/K}\\
$h_{fg}$ & & \num{2.35e6} & \si{J/kg}\\\midrule
\multicolumn{4}{l}{$\sigma=\SI{0.065}{N/m}$, $Pr_t=1.0$, laminar}\\
\bottomrule
\end{tabular}\end{minipage}\hfill
\begin{minipage}[t]{0.42\textwidth}\centering\caption{Run control.}\label{of-tab:run}\small
\begin{tabular}{ll}
\toprule
Setting & Value\\\midrule
$\Delta t$ & \SI{1e-6}{s}\\
$\mathrm{Co},\mathrm{Co}_\alpha$ & $\le0.4$\\
end (demo) & \SI{8e-4}{s}\\
$\alpha$ & MULES, $c_\alpha=1$\\
$p_{rgh}$ & PCG/DIC\\
PISO & 2 correctors\\
$C_C,C_E$ & 1.0,\,0.1 (uncal.)\\
\bottomrule
\end{tabular}\end{minipage}
\end{table}

\section{Reference data and the validation target}
\label{of-sec:reference}
This section assembles what the CFD must reproduce. \textbf{The curves and resistances
here are measured data and the validated reduced-order model of the same apparatus, not
CFD output.} They define the envelope against which the converged CFD
(Section~\ref{of-sec:post}) will be judged.

\subsection{Measured boiling curves and the reduced-order model}
Figure~\ref{of-fig:boil} overlays the measured boiling curves (points) with the reduced-order
model (lines) for both condensers and both chips across coolant set points. The model
tracks the data with a per-configuration RMS error of \SIrange{3.1}{5.2}{K} in surface
temperature. The microchannel surface boils at markedly lower superheat than the plain
surface at the same flux, and the 33-tube condenser reaches higher flux at a given surface
temperature than the confined 42-tube. These are the trends, and the absolute curves,
the 3D simulation is required to recover.

\begin{figure}[H]
\centering\includegraphics[width=0.95\textwidth]{fig_boiling_overlay.png}
\caption{Boiling curves: reduced-order model (lines) vs experiment (points) for both
condensers and both chips, colored by coolant inlet temperature. \emph{Reference data and
validated reduced-order model; not CFD output.}}
\label{of-fig:boil}
\end{figure}

\subsection{Thermal resistance}
The chip-to-coolant path is summarized by the total resistance
$R_{tot}=(T_{surf}-T_{in})/Q$, split into a boiling (chip-side) resistance
$R_{chip}=(T_{surf}-T_{liq})/Q$ and a pool-to-coolant (condenser-side) resistance
$R_{cond}=(T_{liq}-T_{in})/Q$. Figure~\ref{of-fig:res} shows $R_{tot}$ versus heat flux from
experiment and from the reduced-order model; Table~\ref{of-tab:res} lists the median
breakdown. Two results stand out. First, the condenser-side resistance is far smaller for
the open 33-tube bundle (\SIrange{14}{15}{mK/W}) than for the confined 42-tube bundle
(\SIrange{41}{74}{mK/W}), consistent with the open bundle being coolant-side limited and
the confined bundle carrying a layout penalty. Second, the microchannel surface roughly
halves the chip-side resistance relative to plain (e.g. 17 vs \SI{64}{mK/W} on the
33-tube), which is the resistance signature of its lower boiling superheat. The
reduced-order model reproduces $R_{tot}$ to within a few \si{mK/W} except for the 33-tube
microchannel, where it over-predicts (44 vs \SI{33}{mK/W}), the same configuration that
carries the largest boiling-curve error.

\begin{figure}[H]
\centering\includegraphics[width=0.92\textwidth]{fig_resistance.png}
\caption{Total chip-to-coolant resistance $R_{tot}=(T_{surf}-T_{in})/Q$ vs heat flux:
experiment (markers) and reduced-order model (dashed), per condenser.
\emph{Not CFD output.}}
\label{of-fig:res}
\end{figure}

\begin{table}[H]
\centering\caption{Median thermal resistance breakdown (\si{mK/W}). $R_{tot}$ is shown for
experiment and the reduced-order model; the experimental split into $R_{chip}$ and
$R_{cond}$ is from $(T_{surf}-T_{liq})$ and $(T_{liq}-T_{in})$.}
\label{of-tab:res}
\begin{tabular}{lcccc}
\toprule
Configuration & $R_{tot}$ exp & $R_{tot}$ model & $R_{chip}$ exp & $R_{cond}$ exp\\
\midrule
33-tube, plain        & 78.9  & 80.5  & 64.4 & 13.7\\
33-tube, microchannel & 33.1  & 44.1  & 17.3 & 15.1\\
42-tube, plain        & 148.2 & 161.9 & 75.1 & 74.0\\
42-tube, microchannel & 77.1  & 78.7  & 31.9 & 41.1\\
\bottomrule
\end{tabular}
\end{table}

\subsection{Calibration workflow and CHF}
The constant-coefficient closure carries two coefficients that must be calibrated against
data (Fig.~\ref{of-fig:workflow}). Calibration sweeps the chip wall temperature for each
candidate $(C_E,C_C)$, builds the simulated boiling curve from a chip-patch heat-flux
probe ($k_{\mathrm{eff}}=\alpha k_w+(1-\alpha)k_v$), and minimizes the RMS misfit to the
measured curve plus a CHF turn-over penalty. The 42-tube/plain target is 64 measured
points across four coolant temperatures with a measured CHF of \SI{65}{W/cm^2}
(Fig.~\ref{of-fig:val}).

\begin{figure}[H]
\centering
\begin{subfigure}{0.62\textwidth}\includegraphics[width=\textwidth]{fig_workflow.png}
\caption{Refine $\rightarrow$ calibrate $\rightarrow$ run $\rightarrow$ validate.}
\end{subfigure}\hfill
\begin{subfigure}{0.36\textwidth}\includegraphics[width=\textwidth]{fig_validation.png}
\caption{42-tube/plain target.}\end{subfigure}
\caption{Calibration-to-validation workflow and the measured 42-tube/plain calibration
target (CHF \SI{65}{W/cm^2}).}
\label{of-fig:workflow}\label{of-fig:val}
\end{figure}

\section{Condenser coolant-side flow routing and thermal-hydraulics}
\label{of-sec:coolant}
The two condensers route the single-phase coolant differently, which shapes the cold-side
boundary that the fixed-temperature tube walls idealize. The 33-tube bundle uses an open
manifold with a single parallel pass: coolant enters a header spanning the bundle,
distributes across all tubes, flows along the tube axis, and collects at the outlet header.
The 42-tube bundle uses a confined manifold: coolant enters from the top, descends into the
inlet manifold, passes through the tubes along their length, and returns up to the top
through the outlet manifold (Figs.~\ref{of-fig:cooltop} and~\ref{of-fig:coolside}).

The figures show the coolant temperature along the flow path, computed from a
one-dimensional energy balance on the measured condenser duty (Table~\ref{of-tab:coolant}):
the coolant warms from the measured inlet to outlet temperature as it absorbs the condenser
heat load, shown here at the \SI{20}{\celsius}, plain-chip operating point so the two
bundles share an inlet. \textbf{These are coolant-side thermal-hydraulic estimates on the
real geometry, not resolved CFD fields}; in the VOF case the tubes are fixed-temperature
walls and the coolant interior is not simulated. The temperature rise is modest in a single
pass (\SI{1.7}{K} for the 33-tube, \SI{2.1}{K} for the 42-tube at this point, the latter
warmer at outlet because its flow is lower, Table~\ref{of-tab:coolant}) and scales with heat
load, reaching about \SI{5}{K} for the 42-tube microchannel near CHF.

\begin{figure}[H]
\centering\includegraphics[width=0.96\textwidth]{fig_coolant_topview.png}
\caption{Longitudinal top view: coolant temperature along the flow path. The 33-tube (open
manifold) makes a single parallel pass; the 42-tube (confined manifold) enters from the top,
runs through the tubes, and returns up (side $y$--$z$ inset). \emph{Coolant-side
thermal-hydraulic estimate from the measured duty; not a resolved CFD field.}}
\label{of-fig:cooltop}
\end{figure}

The top view (Fig.~\ref{of-fig:cooltop}) shows the longitudinal warming along the tube axis. The
side view (Fig.~\ref{of-fig:coolside}, left column) contrasts the in-plane single pass of the
open 33-tube bundle with the down-through-up routing of the confined 42-tube. The end
cross-sections (Fig.~\ref{of-fig:coolside}, right column) show the bundle packing, the real
four-row, 21-column staggered 42-tube layout and the schematic three-row 33-tube, with the
inlet-to-outlet temperature difference. Resolving the coolant-side velocity, pressure drop,
and header maldistribution would require a separate internal-flow or conjugate simulation
(Section~\ref{of-sec:next}).

\begin{figure}[H]
\centering\includegraphics[width=0.98\textwidth]{fig_coolant_sideend.png}
\caption{Side ($y$--$z$) and end ($x$--$z$) views of the coolant path. Left: the 33-tube
horizontal single pass and the 42-tube down-from-top / through-tubes / back-up routing.
Right: bundle cross-sections at the inlet and outlet ends, with the real staggered 42-tube
layout. \emph{Coolant-side thermal-hydraulic estimate; not a resolved CFD field.}}
\label{of-fig:coolside}
\end{figure}

\section{Post-processing, diagnostics, and field contours}
\label{of-sec:post}
The case is analyzed with a PyVista post-processor that reads the OpenFOAM time
directories through the same reader ParaView uses (Fig.~\ref{of-fig:pipe}). Two data streams
are processed: the instantaneous fields, for interface dynamics and snapshots, and the
\texttt{fieldAverage} (time-averaged) fields, for the stationary quantities that are
compared with data. Geometry and the interface are drawn from the mesh patches and the
$\alpha=0.5$ isosurface.

\begin{figure}[H]
\centering\includegraphics[width=0.95\textwidth]{fig_pipeline.png}
\caption{Post-processing and diagnostics pipeline, from the case time directories to
contours, global balances, and the boiling curve / resistance overlaid on experiment.}
\label{of-fig:pipe}
\end{figure}

\subsection{Chip heat-flux probe and the simulated boiling curve}
\label{of-sec:probe}
The simulated heat flux is the phase-weighted conductive flux integrated over the chip
patch and normalized by its area,
\begin{equation}
q''_{chip}=\frac{1}{A_{chip}}\int_{A_{chip}} k_{\mathrm{eff}}\,\nabla T\!\cdot\!\mathbf n
\,\mathrm dA,
\qquad k_{\mathrm{eff}}=\alpha k_l+(1-\alpha)k_v,
\label{of-eq:probe}
\end{equation}
evaluated on the time-averaged field. Sweeping the chip wall temperature and recording the
stationary $q''_{chip}$ traces out the simulated boiling curve, which is the quantity
fitted to the measured curve during calibration and overlaid on Fig.~\ref{of-fig:boil} for
validation. The boiling heat-transfer coefficient and the simulated thermal resistance
follow directly,
\begin{equation}
h=\frac{q''_{chip}}{T_{surf}-T_{sat}},\qquad
R_{tot}^{\mathrm{CFD}}=\frac{T_{surf}-T_{in}}{q''_{chip}A_{chip}},
\end{equation}
the latter compared against the measured and reduced-order values of Table~\ref{of-tab:res}.

\subsection{Global balances and verification}
At statistical stationarity the chamber must close two balances (Fig.~\ref{of-fig:bal}),
which serve as solver-level verification independent of any experimental comparison: net
evaporation equals net condensation, and the chip heat input equals the condenser heat
removal, both carried by the latent flux,
\begin{equation}
\int_V \dot m_e\,\mathrm dV=\int_V \dot m_c\,\mathrm dV,
\qquad
Q_{chip}=Q_{cond}=h_{fg}\!\int_V \dot m_e\,\mathrm dV .
\end{equation}
The condenser side is cross-checked against the measured coolant calorimetry of
Table~\ref{of-tab:coolant}, $Q_{cond}=\dot m_{cool}c_p(T_{out}-T_{in})$; a converged run that
violates these balances beyond a small tolerance is rejected regardless of how well a
single field matches.

\begin{figure}[H]
\centering\includegraphics[width=0.62\textwidth]{fig_balance.png}
\caption{Steady-state verification balances: evaporation at the chip equals condensation at
the tubes, and the chip heat input equals the condenser duty set by the measured coolant
temperature rise.}
\label{of-fig:bal}
\end{figure}

\subsection{Stationarity, convergence, and mesh sensitivity}
A boiling simulation never reaches a steady state pointwise; the diagnostics are declared
converged when the running means of three integral monitors, the chip heat flux
$q''_{chip}$, the total vapor volume $\int(1-\alpha)\,\mathrm dV$, and the interfacial area
$\int|\nabla\alpha|\,\mathrm dV$, settle within a fixed band over a sliding window, with
their time averages no longer drifting. Mesh sensitivity is then established by repeating
the converged run on the mesh levels of Table~\ref{of-tab:mesh} and reporting the grid
convergence index (GCI) on $q''_{chip}$ and CHF; the production grid is selected where the
GCI falls below the measurement-uncertainty band (Section~\ref{of-sec:next}, item~1).

\subsection{Sampling planes, lines, and derived fields}
Scalar fields are extracted on two planes and two lines (Fig.~\ref{of-fig:planes}): a
horizontal $x$--$y$ plane through the bundle (\emph{top view}) and the vertical $x$--$z$
mid-plane (\emph{side view}) for contours, a vertical line above the chip center for
temperature and velocity profiles through the thermal boundary layer and interface, and a
horizontal midline through the bundle for the tube-to-tube field variation. From these,
the derived fields reported are the heat-transfer coefficient, the interfacial area density
$|\nabla\alpha|$, and the local vapor generation rate $\dot m_e$. The plane locations and
the color treatment are fixed here so the contour figures populate unchanged once the
converged run is available.

\begin{figure}[H]
\centering\includegraphics[width=0.62\textwidth]{fig_sampling.png}
\caption{Sampling geometry: top-view ($x$--$y$, green) and side-view ($x$--$z$, orange)
contour planes, with a vertical profile line above the chip (purple) and the bundle midline
(teal).}
\label{of-fig:planes}
\end{figure}

\subsection{Verification and single-phase reference fields}
Before any phase-change field is trusted, the setup is checked against two fields that are
known independently of the VOF solution (Fig.~\ref{of-fig:ref}). The hydrostatic pressure
$p=\rho_l g\,(z_{fill}-z)$ in the pool is analytic and is the rest-state limit the
$p_{rgh}$ solution must recover with $\mathbf U\to 0$. The single-phase conduction field,
$\nabla\!\cdot(k\nabla T)=0$ with the chip and tubes held at their wall temperatures and
the outer walls adiabatic, is solved here on the real $x$--$z$ geometry and bounds the
thermal field in the no-flow, no-boiling limit; it localizes the conductive resistance
around the chip and the cold sink at the tubes. \textbf{These are verification references,
not the VOF phase-change result}, boiling and convection make the actual pool far more
uniform and concentrate the gradients into thin interfacial layers. They confirm the
boundary-condition setup and provide a conduction-only bound; they do not predict the
boiling fields.

\begin{figure}[H]
\centering\includegraphics[width=0.82\textwidth]{fig_reffields.png}
\caption{Verification / reference fields on the $x$--$z$ plane: (a) single-phase steady
conduction with the chip and tubes at their wall temperatures (no boiling or convection);
(b) analytical hydrostatic pressure. \textbf{Neither is the VOF phase-change solution};
they verify the boundary-condition setup and bound the no-flow limit.}
\label{of-fig:ref}
\end{figure}

\subsection{Field contours (pending the converged run)}
The following are produced by the production run (Section~\ref{of-sec:reference} workflow) and
are intentionally left as reserved slots; no surrogate or unconverged fields are shown as
results:
\begin{itemize}
\item \emph{Temperature contours}, top ($x$--$y$) and side ($x$--$z$) views, for the
42-tube and 33-tube bundles.
\item \emph{Pressure contours}, top and side views, for both bundles.
\item \emph{Velocity-magnitude} and interface ($\alpha$) fields on the same planes, with the
profiles along the sampling lines of Fig.~\ref{of-fig:planes}.
\item \emph{CFD-to-experiment overlay}: the simulated boiling curve and thermal resistance
from the time-averaged probe, Eq.~\eqref{of-eq:probe}, overlaid on Figs.~\ref{of-fig:boil} and
\ref{of-fig:res} for both condensers, with the CHF turn-over checked against the measured
value.
\end{itemize}
The demonstration temperature and velocity fields carry non-physical extremes and are
clamped for display only; they are not presented as results.

\section{A reduced-cost two-dimensional demonstration}
\label{of-sec:twod}
The three-dimensional case in the preceding sections resolves the real bundle but is an
HPC-scale solve. To exercise the same phase-change pipeline at desktop cost, the project
also includes a two-dimensional cross-section case (\texttt{openfoam/chamber42\_2d}). It is
a deliberate simplification with a clearly bounded purpose: exercise the same phase-change
pipeline at desktop cost and establish what such a reduced model can and cannot deliver.

The two-dimensional mesh is generated by \texttt{blockMesh} alone, one cell thick in the
out-of-plane direction with empty front and back patches. The chip is a native floor patch,
obtained by splitting the streamwise direction into three blocks (left floor, chip, right
floor) rather than by feature extraction, so no \texttt{snappyHexMesh}, surface mesh, or
patch surgery is required and the mesh is fully reproducible. The condenser is idealised as
the cold top surface at the bundle-top height, held at the coolant temperature; resolving the
forty-two circular tubes in two dimensions would require \texttt{snappyHexMesh}, whose
isotropic refinement subdivides the single out-of-plane cell and breaks the empty-patch
two-dimensional constraint. The resolved tube bundle therefore remains the three-dimensional
domain, and the two-dimensional case carries the condenser only as a cold boundary. The
operating point is set to the measured 42-tube, plain-chip, \SI{20}{\celsius}-coolant
condition, with saturation \SI{322}{\kelvin} and the cold surface at \SI{293}{\kelvin}.

The intended calibration method is the chip heat-flux probe of Section~\ref{of-sec:probe} driving
a Nelder--Mead search over the evaporation and condensation rate coefficients: the chip wall
temperature is swept across the measured range and the time-averaged chip flux is matched to the
measured curve, with the evaporation coefficient the primary boiling-branch knob. Running this
case on a single workstation core under ESI OpenFOAM~v2512 establishes three things about that
method, each of which bears directly on the HPC plan.

First, the pipeline is sound end to end. The \texttt{blockMesh} mesh passes \texttt{checkMesh},
the solver reads the phase-change model and advances stably with the temperature field bounded
between the coolant and chip values, and the coded probe compiles at run time and records the
chip heat flux. The machinery built for the three-dimensional case runs unmodified on the
two-dimensional surrogate.

Second, driven beyond the measured critical heat flux the model reproduces the boiling crisis.
With the chip held \SI{36}{\kelvin} above saturation, past the measured
\SI{65}{\watt\per\centi\meter\squared} CHF, vapour accumulates at the surface, coalesces into a
film after roughly \SI{60}{\milli\second}, and the chip heat flux collapses toward zero as the
low-conductivity vapour blankets the wall. This dryout is the one regime in which the phase
change visibly governs the surface flux.

Third, and decisively for the calibration strategy, on the nucleate branch the chip flux is
conduction-dominated and essentially independent of the Lee coefficient at this mesh resolution.
Holding the wall at \SI{65}{\celsius} (\SI{16}{\kelvin} superheat) and tripling the evaporation
coefficient from \num{0.1} to \num{0.3} leaves the chip heat flux unchanged to four significant
figures and the domain vapour fraction unchanged in the sixth decimal place
(Table~\ref{of-tab:2dsens}). The subcooled pool and the absence of resolved nucleation sites mean
that only a thin near-wall band of cells ever exceeds saturation, so the integrated evaporation
source is negligible and the surface flux is set by transient conduction into the slowly warming
pool rather than by boiling. That bulk warm-up is a seconds-scale diffusion process across the
\SI{24}{\milli\meter} gap, so at the sub-second times affordable on a single core the chip flux
has not stationarised, and no value of the rate coefficient changes that.

\begin{table}[ht]
\centering
\caption{Chip heat flux at a \SI{65}{\celsius} wall (\SI{16}{\kelvin} superheat), fixed
\SI{0.1}{\second} solve. Tripling the Lee evaporation coefficient leaves both the surface flux
and the domain vapour inventory essentially unchanged: at this two-dimensional mesh resolution
the phase-change source has no leverage on the measured quantity.}
\label{of-tab:2dsens}
\begin{tabular}{lcc}
\toprule
Lee evaporation coeff.\ $C_e$ & Chip flux $q''$ (\si{\watt\per\centi\meter\squared}) & Domain liquid fraction \\
\midrule
0.1 & 10.30 & 0.7385951 \\
0.3 & 10.30 & 0.7385937 \\
\bottomrule
\end{tabular}
\end{table}

\begin{figure}[ht]
\centering
\begin{subfigure}{0.49\textwidth}
\includegraphics[width=\linewidth]{twod_nucleate}
\caption{Nucleate point (\SI{65}{\celsius} wall, \SI{16}{\kelvin} superheat): the chip flux
decays monotonically and has not stationarised at \SI{0.1}{\second}.}
\end{subfigure}\hfill
\begin{subfigure}{0.49\textwidth}
\includegraphics[width=\linewidth]{twod_collapse}
\caption{Past CHF (\SI{85}{\celsius} wall, \SI{36}{\kelvin} superheat): a vapour film forms near
\SI{63}{\milli\second} and the chip flux collapses toward zero (boiling crisis).}
\end{subfigure}
\caption{Two-dimensional demonstration runs on a single core under ESI OpenFOAM~v2512. The
nucleate branch is conduction-dominated and slow to settle; driven past the measured CHF the
model exhibits vapour-film dryout. Both traces are early-transient and non-converged, consistent
with the Lee-coefficient insensitivity of Table~\ref{of-tab:2dsens}.}
\label{of-fig:twod}
\end{figure}

The consequence is that the quantitative calibration cannot be performed on the two-dimensional
surrogate (Figure~\ref{of-fig:twod}): the coefficient the calibration must fit has no leverage on the quantity it is fitted
against, because the mesh does not resolve the boiling. This is a sharper statement of the HPC
requirement than cost alone. The resolved-bubble, snapped-and-layered three-dimensional mesh
described above is not merely faster-converging on a cluster; it is the resolution at which the
Lee source becomes a dominant term in the surface energy balance and the calibration becomes
well-posed. The two-dimensional case is therefore a validated demonstration that the pipeline
runs and behaves physically at its limits, and direct evidence for where the quantitative work
must be done, not a substitute for it.

\section{Next steps}
\label{of-sec:next}
\begin{enumerate}
\item \textbf{Quantify the measurement-uncertainty band} (heat flux \SIrange{4}{11}{\%},
surface temperature $\pm\SI{0.2}{K}$, CHF $\sim\pm\SI{4}{\%}$). This is the prerequisite
against which any CFD-to-experiment agreement is defined.
\item \textbf{Refine and run the 42-tube case} on HPC: snapped/layered mesh at
\SIrange{0.1}{0.3}{mm}, calibrate $C_E$/$C_C$ against the 42-tube/plain boiling curve and
CHF, and run to a statistically stationary state.
\item \textbf{Produce the field contours} (temperature, pressure, velocity, interface) on
the top and side planes of Fig.~\ref{of-fig:planes} and insert them in Section~\ref{of-sec:post}.
\item \textbf{Build and run the 33-tube CFD case} (parallel geometry, open manifold) so the
cross-condenser comparison, larger/fewer tubes versus smaller/denser, is resolved in 3D and
not only in the reduced-order model.
\item \textbf{Resolve the coolant side} with a separate internal-flow or conjugate
heat-transfer simulation of each manifold and tube bundle, replacing the fixed-temperature
wall with a computed tube-wall temperature and capturing coolant velocity, pressure drop, and
header maldistribution; the thermal-hydraulic estimate of Section~\ref{of-sec:coolant} is the
one-dimensional limit of this.
\item \textbf{Overlay the converged CFD on experiment} for both condensers
(Figs.~\ref{of-fig:boil}, \ref{of-fig:res}) and report the CFD thermal resistance against
Table~\ref{of-tab:res}; passing requires lying within the uncertainty band point by point.
\item \textbf{Add the microchannel surface} to the CFD, either resolved or as a calibrated
sub-grid enhancement, and validate against the microchannel boiling curves.
\item \textbf{Resolve the CHF mechanism} (true burnout versus a surface-temperature
ceiling) from the wall-temperature traces, and reconcile the 33-tube/66-tube CHF behavior.
\item \textbf{Extend to the 66-tube bundle and HFE-7000} once those datasets are available,
to test cross-geometry and cross-fluid generalization in 3D.
\end{enumerate}

\section{Status and outlook}
\label{of-sec:status}
This document establishes a reproducible 3D VOF phase-change framework on the true
geometry: a verified body-fitted meshing pipeline (\num{299656} cells, \texttt{checkMesh}
OK), a complete physics, geometry, and numerics specification, the measured condenser
coolant operating conditions, and the full set of experimental and reduced-order
validation targets, boiling curves and thermal resistances for both condensers. What
remains is the HPC campaign of Section~\ref{of-sec:next}: calibration, the run to stationarity,
the field contours, and the CFD-to-experiment overlay, together with the independent
prerequisite of quantifying the uncertainty band. On completion, the framework yields
resolved boiling and condensation fields that complement, and stress-test, the validated
reduced-order model of the same apparatus.



\clearpage
\begin{thebibliography}{99}
\small
%-------- Part I references --------
\bibitem{mudawar2001} Mudawar, I., 2001, ``Assessment of high-heat-flux thermal management schemes,'' \textit{IEEE Trans. Components and Packaging Technologies}, 24(2), pp.~122--141.
\bibitem{azarifar2024} Azarifar, M., Arik, M., and Chang, J.-Y., 2024, ``Liquid cooling of data centers: A necessity facing challenges,'' \textit{Applied Thermal Engineering}, 247, 123112.
\bibitem{birbarah2020} Birbarah, P., Gebrael, T., Foulkes, T., et al., 2020, ``Water immersion cooling of high power density electronics,'' \textit{Int.\ J.\ Heat and Mass Transfer}, 147, 118918.
\bibitem{wu2025} Wu, X., et al., 2025, ``Investigations on heat dissipation performance and overall characteristics of two-phase liquid immersion cooling systems for data center,'' \textit{Int.\ J.\ Heat and Mass Transfer}, 239, 126575.
\bibitem{mody2022} Mody, F., Chauhan, A., Shukla, M.~Y., and Kandlikar, S.~G., 2022, ``Evaluation of heater size and external enhancement techniques in pool boiling heat transfer with dielectric fluids,'' \textit{Int.\ J.\ Heat and Mass Transfer}, 191, 122857.
\bibitem{shukla2024} Shukla, M.~Y., and Kandlikar, S.~G., 2024, ``Novel subcooled boiling chamber with submerged condensation for high heat flux removal for data center application,'' \textit{Proc.\ 23rd IEEE Intersociety Conf.\ on Thermal and Thermomechanical Phenomena in Electronic Systems (ITherm)}, Orlando, FL.
\bibitem{shrigondekar2023} Shrigondekar, H., Lin, Y.-C., and Wang, C.-C., 2023, ``Investigations on performance of single-phase immersion cooling system,'' \textit{Int.\ J.\ Heat and Mass Transfer}, 206, 123961.
\bibitem{nukiyama1934} Nukiyama, S., 1934, ``The maximum and minimum values of the heat transmitted from metal to boiling water under atmospheric pressure,'' \textit{J.\ Soc.\ Mech.\ Eng.\ Japan}, 37, pp.~367--374 (reprinted \textit{Int.\ J.\ Heat Mass Transfer}, 9, 1966, pp.~1419--1433).
\bibitem{rohsenow1952} Rohsenow, W.~M., 1952, ``A method of correlating heat transfer data for surface boiling of liquids,'' \textit{Trans.\ ASME}, 74, pp.~969--976.
\bibitem{vachon1968} Vachon, R.~I., Nix, G.~H., and Tanger, G.~E., 1968, ``Evaluation of constants for the Rohsenow pool-boiling correlation,'' \textit{J.\ Heat Transfer}, 90(2), pp.~239--247.
\bibitem{mikic1969} Mikic, B.~B., and Rohsenow, W.~M., 1969, ``A new correlation of pool boiling data including the effect of heating surface characteristics,'' \textit{J.\ Heat Transfer}, 91(2), pp.~245--250.
\bibitem{pioro1999} Pioro, I.~L., 1999, ``Experimental evaluation of constants for the Rohsenow pool boiling correlation,'' \textit{Int.\ J.\ Heat and Mass Transfer}, 42(11), pp.~2003--2013.
\bibitem{carey2008} Carey, V.~P., 2008, \textit{Liquid--Vapor Phase-Change Phenomena}, 2nd ed., Taylor \& Francis, New York.
\bibitem{incropera2007} Incropera, F.~P., DeWitt, D.~P., Bergman, T.~L., and Lavine, A.~S., 2007, \textit{Fundamentals of Heat and Mass Transfer}, 6th ed., Wiley, Hoboken, NJ.
\bibitem{zuber1959} Zuber, N., 1959, ``Hydrodynamic aspects of boiling heat transfer,'' Ph.D.\ dissertation, UCLA / AEC Report AECU-4439.
\bibitem{kutateladze1948} Kutateladze, S.~S., 1948, ``On the transition to film boiling under natural convection,'' \textit{Kotloturbostroenie}, 3, pp.~10--12.
\bibitem{lienhard1973} Lienhard, J.~H., and Dhir, V.~K., 1973, ``Extended hydrodynamic theory of the peak and minimum pool boiling heat fluxes,'' NASA CR-2270.
\bibitem{kandlikar2001} Kandlikar, S.~G., 2001, ``A theoretical model to predict pool boiling CHF incorporating effects of contact angle and orientation,'' \textit{J.\ Heat Transfer}, 123(6), pp.~1071--1079.
\bibitem{mudawar1997} Mudawar, I., Howard, A.~H., and Gersey, C.~O., 1997, ``An analytical model for near-saturated pool boiling critical heat flux on vertical surfaces,'' \textit{Int.\ J.\ Heat and Mass Transfer}, 40(10), pp.~2327--2339.
\bibitem{cooke2011} Cooke, D., and Kandlikar, S.~G., 2011, ``Pool boiling heat transfer and bubble dynamics over plain and enhanced microchannels,'' \textit{J.\ Heat Transfer}, 133(5), 052902.
\bibitem{cooke2012} Cooke, D., and Kandlikar, S.~G., 2012, ``Effect of open microchannel geometry on pool boiling enhancement,'' \textit{Int.\ J.\ Heat and Mass Transfer}, 55(4), pp.~1004--1013.
\bibitem{elgenk2021a} El-Genk, M.~S., and Pourghasemi, M., 2021, ``Experiments and correlations of saturation boiling of HFE-7000 dielectric liquid on rough inclined copper surfaces,'' \textit{Int.\ J.\ Heat and Mass Transfer}, 164, 120540.
\bibitem{elgenk2021b} El-Genk, M.~S., and Pourghasemi, M., 2021, ``Subcooled boiling critical heat flux of HFE-7000 dielectric liquid on inclined rough Cu,'' \textit{Int.\ J.\ Heat and Mass Transfer}, 175, 121354.
\bibitem{jaikumar2015} Jaikumar, A., and Kandlikar, S.~G., 2015, ``Enhanced pool boiling heat transfer mechanisms for selectively sintered open microchannels,'' \textit{Int.\ J.\ Heat and Mass Transfer}, 88, pp.~652--661.
\bibitem{patil2014} Patil, C.~M., and Kandlikar, S.~G., 2014, ``Pool boiling enhancement through microporous coatings selectively electrodeposited on fin tops of open microchannels,'' \textit{Int.\ J.\ Heat and Mass Transfer}, 79, pp.~816--828.
\bibitem{jaikumar2016} Jaikumar, A., and Kandlikar, S.~G., 2016, ``Pool boiling enhancement through bubble induced convective liquid flow in feeder microchannels,'' \textit{Applied Physics Letters}, 108(4), 041604.
\bibitem{nusselt1916} Nusselt, W., 1916, ``Die Oberfl\"achenkondensation des Wasserdampfes,'' \textit{Z.\ Ver.\ Dtsch.\ Ing.}, 60, pp.~541--546, 569--575.
\bibitem{rose2002} Rose, J.~W., 2002, ``Dropwise condensation theory and experiment: A review,'' \textit{Proc.\ IMechE Part A: J.\ Power and Energy}, 216(2), pp.~115--128.
\bibitem{collier1994} Collier, J.~G., and Thome, J.~R., 1994, \textit{Convective Boiling and Condensation}, 3rd ed., Oxford University Press, New York.
\bibitem{hausen1943} Hausen, H., 1943, ``Darstellung des W\"arme\"uberganges in Rohren durch verallgemeinerte Potenzbeziehungen,'' \textit{VDI-Zeitung, Beiheft Verfahrenstechnik}, 4, pp.~91--98.
\bibitem{shah1978} Shah, R.~K., and London, A.~L., 1978, \textit{Laminar Flow Forced Convection in Ducts}, Academic Press, New York.
\bibitem{gnielinski1976} Gnielinski, V., 1976, ``New equations for heat and mass transfer in turbulent pipe and channel flow,'' \textit{Int.\ Chemical Engineering}, 16(2), pp.~359--368.
\bibitem{churchill1975} Churchill, S.~W., and Chu, H.~H.~S., 1975, ``Correlating equations for laminar and turbulent free convection from a vertical plate,'' \textit{Int.\ J.\ Heat and Mass Transfer}, 18(11), pp.~1323--1329.
\bibitem{faghri1995} Faghri, A., 1995, \textit{Heat Pipe Science and Technology}, Taylor \& Francis, Washington, DC.

%-------- Part II references --------
\bibitem{of-shukla2024} Shukla, M.~Y., and Kandlikar, S.~G., 2024, ``Novel subcooled boiling
chamber with submerged condensation for high heat flux removal for data center
application,'' \textit{Proc.\ 23rd IEEE ITherm}.
\bibitem{of-weller1998} Weller, H.~G., Tabor, G., Jasak, H., and Fureby, C., 1998, ``A
tensorial approach to computational continuum mechanics using object-oriented
techniques,'' \textit{Comput.\ Phys.}, 12(6), pp.~620--631.
\bibitem{of-hirt1981} Hirt, C.~W., and Nichols, B.~D., 1981, ``Volume of fluid (VOF) method
for the dynamics of free boundaries,'' \textit{J.\ Comput.\ Phys.}, 39(1), pp.~201--225.
\bibitem{of-brackbill1992} Brackbill, J.~U., Kothe, D.~B., and Zemach, C., 1992, ``A continuum
method for modeling surface tension,'' \textit{J.\ Comput.\ Phys.}, 100(2), pp.~335--354.
\bibitem{of-lee1980} Lee, W.~H., 1980, ``A pressure iteration scheme for two-phase flow
modeling,'' in \textit{Multiphase Transport}, Hemisphere.
\bibitem{of-cooke2012} Cooke, D., and Kandlikar, S.~G., 2012, ``Effect of open microchannel
geometry on pool boiling enhancement,'' \textit{Int.\ J.\ Heat Mass Transfer}, 55(4),
pp.~1004--1013.
\bibitem{of-kandlikar2001} Kandlikar, S.~G., 2001, ``A theoretical model to predict pool
boiling CHF incorporating effects of contact angle and orientation,'' \textit{J.\ Heat
Transfer}, 123(6), pp.~1071--1079.
\bibitem{of-mody2022} Mody, F., Chauhan, A., Shukla, M.~Y., and Kandlikar, S.~G., 2022,
``Evaluation of heater size and external enhancement techniques in pool boiling heat
transfer with dielectric fluids,'' \textit{Int.\ J.\ Heat Mass Transfer}, 191, 122857.

\end{thebibliography}

\end{document}
